% Les pétroles — PCT 3ème — Cours signé OnBuch+
% Programme officiel MINESEC. Compiler avec : tectonic N16-les-petroles.tex
\newcommand{\DOCMATIERE}{PCT}
\newcommand{\DOCNIVEAU}{3ème}
\newcommand{\DOCMODULE}{Module 3 : Chimie appliquée}
\newcommand{\DOCLECON}{N16}
\newcommand{\DOCTITRE}{Les pétroles}
% OnBuch+ — préambule « cours signés » ENRICHI (compilé avec tectonic / XeLaTeX)
% Système visuel « Pop » d'OnBuch+ : fond crème (jamais blanc), bordures encre
% épaisses, ombres « dures » décalées, titres Archivo Black en capitales.
% Version MATHÉMATIQUES : graphes (pgfplots/tikz), boîtes pédagogiques
% (définition, propriété, méthode, attention, le savais-tu, exercice résolu,
% exercice, TP, bilan), page de garde et signature OnBuch+ ancrée en bas.
%
% Macros attendues dans main.tex AVANT \input{preamble} :
%   \DOCMATIERE  \DOCNIVEAU  \DOCMODULE  \DOCLECON (numéro)  \DOCTITRE
\documentclass[11pt]{article}

\usepackage{fontspec}
\usepackage[french]{babel}
\usepackage[a4paper,margin=2cm,top=3.4cm,bottom=2.8cm,headheight=1.4cm,footskip=1.3cm]{geometry}
\usepackage{amsmath,amssymb,amsfonts}
\usepackage{mathtools} % \xmapsto, \coloneqq... souvent utilisés par les modèles
\usepackage{cancel} % simplifications barrées (bilans de forces)
% \abs{z} (module, valeur absolue) et \norm{u} (norme), utilisés par les modèles
\providecommand{\abs}{}\let\abs\relax\DeclarePairedDelimiter{\abs}{\lvert}{\rvert}
\providecommand{\norm}{}\let\norm\relax\DeclarePairedDelimiter{\norm}{\lVert}{\rVert}
\usepackage[table]{xcolor} % [table] : fournit aussi \rowcolors (alternance de lignes)
\usepackage{pagecolor}
\usepackage{tikz}
\usetikzlibrary{arrows.meta,positioning,calc,shapes.geometric,shapes.misc,shapes.arrows,decorations.pathmorphing,decorations.pathreplacing,decorations.markings,patterns,shadows,mindmap,trees,fit,backgrounds,matrix,intersections,angles,quotes}
\usepackage{pgfplots}
\usepackage[version=4]{mhchem} % équations nucléaires \ce{^{235}_{92}U}
\usepackage[european,siunitx]{circuitikz} % schémas électriques
\pgfplotsset{compat=1.18}
\usepgfplotslibrary{fillbetween,groupplots} % aires entre courbes (calcul intégral)
\usepackage{enumitem}
\usepackage{fancyhdr}
% [most] charge toutes les bibliothèques tcolorbox (listings, minted, raster,
% documentation...) et gonfle inutilement la mémoire TeX sur un document long
% (~300 boîtes) : on ne charge que ce qui est réellement utilisé.
\usepackage[skins,breakable]{tcolorbox}
\usepackage{graphicx}
\usepackage{colortbl}
\usepackage{booktabs}
\usepackage{tabularx}
\usepackage{multirow}
\usepackage{diagbox} % case d'en-tête diagonale des tableaux à double entrée
% Colonnes extensibles centrée / gauche / droite (C, L, R), souvent utilisées par les modèles
\newcolumntype{C}{>{\centering\arraybackslash}X}
\newcolumntype{L}{>{\raggedright\arraybackslash}X}
\newcolumntype{R}{>{\raggedleft\arraybackslash}X}
\usepackage{array}
\usepackage{multicol}
\usepackage{siunitx}
% parse-numbers=false : accepte aussi \SI{2,5 \times 10^{-3}}{...}
\sisetup{locale=FR,output-decimal-marker={,},group-minimum-digits=4,parse-numbers=false}
\DeclareSIUnit\ohm{\ensuremath{\symup{\Omega}}} % U+2126 absent de la police
\DeclareSIUnit\division{div} % réglages d'oscilloscope (ms/div, V/div)
\DeclareSIUnit\tr{tr} % tours (vitesse de rotation, ex. \SI{1500}{\tr\per\minute})
\DeclareSIUnit\molar{M} % concentration molaire (ex. \SI{3}{\molar}, \SI{1}{\micro\molar})
\DeclareSIUnit\an{an} % année (le modèle confond parfois avec \year, un compteur TeX) — ex. \SI{5}{\kilo\meter\per\an}
\DeclareSIUnit\FCFA{FCFA} % franc CFA (ex. \SI{500}{\FCFA\per\kilo\gram})
\DeclareSIUnit\degreeBrix{\textdegree Brix} % teneur en sucre d'un sirop/jus (réfractométrie)
\DeclareSIUnit\month{mois} % le modèle utilise parfois \month, un compteur TeX, comme unité de durée
\DeclareSIUnit\microliter{\micro\liter} % le modèle écrit parfois \microliter en un seul token
\DeclareSIUnit\million{millions} % dénombrement (ex. \SI{15}{\million\per\mL} spermatozoïdes)
\usepackage{qrcode}
% \degree : commande fréquemment utilisée par les modèles pour un angle ou une
% température en dehors de \SI{}{} (non fournie par les paquets chargés).
\providecommand{\degree}{^{\circ}}
% \qed (fin de démonstration), utilisé par les modèles sans amsthm
\providecommand{\qed}{\hfill$\square$}
% \barre : double barre des valeurs interdites dans un tableau de variations (tkz-tab)
\providecommand{\barre}{\ensuremath{\|}}
% environnement proof (amsthm non chargé) utilisé par les modèles
\ifdefined\proof\else\newenvironment{proof}[1][Démonstration]{\par\noindent\textit{#1.}\ }{\qed\par}\fi
\usepackage{lastpage}
\usepackage{hyperref}
\hypersetup{hidelinks}

% ============================================================
% Palette « Pop » OnBuch+ — mêmes hex que la classe P (plus_ui.dart)
% ============================================================
\definecolor{poporange}{HTML}{F59321}
\definecolor{popdark}{HTML}{E07A0C}
\definecolor{popcream}{HTML}{FFF6E8}
\definecolor{popink}{HTML}{1C1714}
\definecolor{poppurple}{HTML}{7A5AE0}
\definecolor{popgreen}{HTML}{2E9E6A}
\definecolor{poppink}{HTML}{E0457B}
\definecolor{popblue}{HTML}{2F7DE1}
\definecolor{popgold}{HTML}{C98A12}
\definecolor{popmuted}{HTML}{8A7F76}
\definecolor{popline}{HTML}{EADFD2}
\definecolor{popgray}{HTML}{8A7F76}
\colorlet{popgrey}{popgray}
% Alias tolérants (le modèle nomme parfois les couleurs différemment)
\colorlet{popwhite}{white}
\colorlet{popblack}{popink}
\colorlet{popred}{poppink}
\colorlet{popyellow}{popgold}
% Teintes claires (fonds de tableaux/encadrés) — le modèle nomme parfois
% les couleurs avec un suffixe L (ex. popblueL) sans les avoir définies.
\colorlet{poporangeL}{poporange!20}
\colorlet{popdarkL}{popdark!20}
\colorlet{poppurpleL}{poppurple!20}
\colorlet{popgreenL}{popgreen!20}
\colorlet{poppinkL}{poppink!20}
\colorlet{popblueL}{popblue!20}
\colorlet{popgoldL}{popgold!20}
\colorlet{popredL}{popred!20}
\colorlet{popgrayL}{popgray!20}
% Teintes claires pour fonds de tableaux / zones
\colorlet{popblueL}{popblue!10!white}
\colorlet{poporangeL}{poporange!14!white}
\colorlet{popgreenL}{popgreen!12!white}
\colorlet{poppurpleL}{poppurple!12!white}
\colorlet{poppinkL}{poppink!12!white}
\colorlet{popgoldL}{popgold!14!white}

\pagecolor{popcream}

% ============================================================
% Typographie
% ============================================================
\defaultfontfeatures{Ligatures=TeX}
\setmainfont{PlusJakartaSans-Medium.ttf}[Path=../../../fonts/,BoldFont=PlusJakartaSans-Bold.ttf]
% Maths en sans-serif (Fira Math) avec chiffres et lettres droites de Plus
% Jakarta, pour que les formules se fondent dans le texte.
\usepackage[math-style=ISO,bold-style=ISO]{unicode-math}
\setmathfont{FiraMath-Regular.otf}[Path=../../../fonts/]
\setmathfont{PlusJakartaSans-Medium.ttf}[Path=../../../fonts/,range={up/num,up/{Latin,latin},\mathup}]
\usepackage{icomma} % virgule décimale sans espace en mode math
% Caractères mathématiques Unicode absents des polices de texte (Plus Jakarta,
% Archivo Black) : rendus par leur équivalent mathématique, en texte comme en
% formule — sinon ils s'affichent en carré vide (ex. « Arithmétique dans ℤ »).
\usepackage{newunicodechar}
\newunicodechar{ℝ}{\ensuremath{\mathbb{R}}}
\newunicodechar{ℤ}{\ensuremath{\mathbb{Z}}}
\newunicodechar{ℂ}{\ensuremath{\mathbb{C}}}
\newunicodechar{ℕ}{\ensuremath{\mathbb{N}}}
\newunicodechar{ℚ}{\ensuremath{\mathbb{Q}}}
\newunicodechar{𝔻}{\ensuremath{\mathbb{D}}}
\newunicodechar{α}{\ensuremath{\alpha}}
\newunicodechar{β}{\ensuremath{\beta}}
\newunicodechar{γ}{\ensuremath{\gamma}}
\newunicodechar{δ}{\ensuremath{\delta}}
\newunicodechar{ε}{\ensuremath{\varepsilon}}
\newunicodechar{θ}{\ensuremath{\theta}}
\newunicodechar{λ}{\ensuremath{\lambda}}
\newunicodechar{μ}{\ensuremath{\mu}}
\newunicodechar{σ}{\ensuremath{\sigma}}
\newunicodechar{τ}{\ensuremath{\tau}}
\newunicodechar{φ}{\ensuremath{\varphi}}
\newunicodechar{ψ}{\ensuremath{\psi}}
\newunicodechar{ω}{\ensuremath{\omega}}
\newunicodechar{Σ}{\ensuremath{\Sigma}}
% majuscules produites par \MakeUppercase dans les titres (α-aminés -> Α)
\newunicodechar{Α}{\ensuremath{\alpha}}
\newunicodechar{Β}{\ensuremath{\beta}}
\newunicodechar{Γ}{\ensuremath{\Gamma}}
\newunicodechar{Δ}{\ensuremath{\Delta}}
\newunicodechar{Ε}{\ensuremath{\varepsilon}}
\newunicodechar{Θ}{\ensuremath{\Theta}}
\newunicodechar{Λ}{\ensuremath{\Lambda}}
\newunicodechar{Μ}{\ensuremath{\mu}}
\newunicodechar{Τ}{\ensuremath{\tau}}
\newunicodechar{Φ}{\ensuremath{\Phi}}
\newunicodechar{Ψ}{\ensuremath{\Psi}}
\newunicodechar{Ω}{\ensuremath{\Omega}}
\newunicodechar{↦}{\ensuremath{\mapsto}}
\newunicodechar{∈}{\ensuremath{\in}}
\newunicodechar{∉}{\ensuremath{\notin}}
\newunicodechar{∧}{\ensuremath{\wedge}}
\newunicodechar{⇔}{\ensuremath{\Leftrightarrow}}
\newunicodechar{⟺}{\ensuremath{\Longleftrightarrow}}
\newunicodechar{⇒}{\ensuremath{\Rightarrow}}
\newunicodechar{⟹}{\ensuremath{\Longrightarrow}}
\newunicodechar{∘}{\ensuremath{\circ}}
\newunicodechar{∪}{\ensuremath{\cup}}
\newunicodechar{∩}{\ensuremath{\cap}}
\newunicodechar{≡}{\ensuremath{\equiv}}
\newunicodechar{⊂}{\ensuremath{\subset}}
\newunicodechar{⊥}{\ensuremath{\perp}}
\newunicodechar{∃}{\ensuremath{\exists}}
\newunicodechar{∀}{\ensuremath{\forall}}
\newunicodechar{∛}{\ensuremath{\sqrt[3]{\,}}}
\newunicodechar{∜}{\ensuremath{\sqrt[4]{\,}}}
\newunicodechar{⃗}{\ensuremath{{}^{\to}}}
\newunicodechar{ⁿ}{\ifmmode^{n}\else\textsuperscript{\ensuremath{n}}\fi}
\newunicodechar{ˣ}{\ifmmode^{x}\else\textsuperscript{\ensuremath{x}}\fi}
\newunicodechar{ᵇ}{\ifmmode^{b}\else\textsuperscript{\ensuremath{b}}\fi}
\newunicodechar{ʳ}{\ifmmode^{r}\else\textsuperscript{\ensuremath{r}}\fi}
\newunicodechar{ᵘ}{\ifmmode^{u}\else\textsuperscript{\ensuremath{u}}\fi}
\newunicodechar{ʸ}{\ifmmode^{y}\else\textsuperscript{\ensuremath{y}}\fi}
\newunicodechar{ᵃ}{\ifmmode^{a}\else\textsuperscript{\ensuremath{a}}\fi}
\newunicodechar{ᵉ}{\ifmmode^{e}\else\textsuperscript{\ensuremath{e}}\fi}
\newunicodechar{ᵏ}{\ifmmode^{k}\else\textsuperscript{\ensuremath{k}}\fi}
\newunicodechar{ᵖ}{\ifmmode^{p}\else\textsuperscript{\ensuremath{p}}\fi}
\newunicodechar{ᴺ}{\ifmmode^{N}\else\textsuperscript{\ensuremath{N}}\fi}
\newunicodechar{ᶜ}{\ifmmode^{c}\else\textsuperscript{\ensuremath{c}}\fi}
\newunicodechar{ʰ}{\ifmmode^{h}\else\textsuperscript{\ensuremath{h}}\fi}
\newunicodechar{ᵠ}{\ifmmode^{q}\else\textsuperscript{\ensuremath{q}}\fi}
\newunicodechar{ₙ}{\ifmmode_{n}\else\textsubscript{\ensuremath{n}}\fi}
\newunicodechar{ₐ}{\ifmmode_{a}\else\textsubscript{\ensuremath{a}}\fi}
\newunicodechar{ᵢ}{\ifmmode_{i}\else\textsubscript{\ensuremath{i}}\fi}
\newunicodechar{ᵣ}{\ifmmode_{r}\else\textsubscript{\ensuremath{r}}\fi}
\newunicodechar{ₖ}{\ifmmode_{k}\else\textsubscript{\ensuremath{k}}\fi}
\newunicodechar{ᵦ}{\ifmmode_{\beta}\else\textsubscript{\ensuremath{\beta}}\fi}
\newunicodechar{ₓ}{\ifmmode_{x}\else\textsubscript{\ensuremath{x}}\fi}
\newfontfamily\popdisplay{ArchivoBlack-Regular.ttf}[Path=../../../fonts/]
\newfontfamily\poplogo{SpaceGrotesk-Bold.ttf}[Path=../../../fonts/]
\newfontfamily\popsemi{PlusJakartaSans-SemiBold.ttf}[Path=../../../fonts/]
\newfontfamily\popxbold{PlusJakartaSans-ExtraBold.ttf}[Path=../../../fonts/]
\newfontfamily\popxboldls{PlusJakartaSans-ExtraBold.ttf}[Path=../../../fonts/,LetterSpace=3.0]

\setlist[enumerate]{leftmargin=1.7em,itemsep=5pt,topsep=4pt}
\setlist[itemize]{leftmargin=1.7em,itemsep=4pt,topsep=4pt,label={\color{poporange}\textbullet}}
\setlength{\parindent}{0pt}
\setlength{\parskip}{5pt}
\renewcommand{\arraystretch}{1.25}

% pgfplots : style Pop par défaut
\pgfplotsset{
  every axis/.append style={axis line style={popink,line width=1pt},tick style={popink},
    grid style={popline,line width=0.5pt},label style={font=\small},tick label style={font=\footnotesize}},
  popcurve/.style={poporange,line width=1.8pt,no markers,smooth},
}
% Même style disponible dans un \draw TikZ simple (hors pgfplots)
\tikzset{popcurve/.style={poporange,line width=1.8pt}}
\tikzset{popgrid/.style={popline,very thin}}
\colorlet{popcurve}{poporange} % le nom du style est parfois utilisé comme couleur

% ============================================================
% Pill / squiggle
% ============================================================
\newcommand{\poppill}[3]{%
  \tikz[baseline=-0.55ex]\node[draw=popink,line width=1.2pt,rounded corners=2.6pt,%
    fill=#2,inner xsep=6.5pt,inner ysep=3pt]{%
    \popxboldls\fontsize{7}{7}\selectfont\color{#3}\MakeUppercase{#1}};%
}
\newcommand{\popsquiggle}[1]{%
  \tikz[baseline=-0.4ex]{
    \draw[#1,line width=2.6pt,line cap=round]
      plot [smooth] coordinates {(0,0.09) (0.34,0.01) (0.68,0.21) (1.02,0.01) (1.36,0.10)};
  }%
}
% Mot-clé surligné (marqueur orange sous le texte)
\newcommand{\cle}[1]{{\popxbold\color{popdark}#1}}

% ============================================================
% En-tête / pied
% ============================================================
\pagestyle{fancy}
\fancyhf{}
\renewcommand{\headrulewidth}{0pt}
\renewcommand{\footrulewidth}{0pt}
\lhead{\raisebox{-0.15ex}{\poplogo\fontsize{13}{13}\selectfont\color{popink}OnBuch\color{poporange}+}}
\rhead{\poppill{\DOCMATIERE\ \textbullet\ \DOCNIVEAU\ \textbullet\ Leçon \DOCLECON}{white}{popink}}
\lfoot{\popxbold\fontsize{7.6}{7.6}\selectfont\color{popmuted}\MakeUppercase{Cours signé OnBuch+}\ \textbullet\ {\popsemi\DOCTITRE}}
\rfoot{\tikz[baseline=-0.6ex]\node[rounded corners=8.5pt,fill=popink,minimum height=17pt,minimum width=17pt,inner xsep=5pt,inner sep=0]{\popxbold\fontsize{8}{8}\selectfont\color{popcream}\thepage\,/\,\pageref*{LastPage}};}

% ============================================================
% Titres
% ============================================================
\newcommand{\chaptitre}[2]{%
  \begin{tcolorbox}[enhanced,colback=poporange,colframe=popink,boxrule=2.6pt,arc=11pt,
      drop shadow=popink,left=15pt,right=15pt,top=12pt,bottom=11pt,before skip=3pt,after skip=18pt]
    \poppill{#2}{popink}{popcream}\par\vspace{9pt}
    {\raggedright\hyphenpenalty=10000\exhyphenpenalty=10000\popdisplay\fontsize{22}{24}\selectfont\color{popcream}\MakeUppercase{#1}\par}
  \end{tcolorbox}
}
% Section numérotée : pastille orange avec le numéro + titre Archivo
\newcounter{popsec}
\newcommand{\coursec}[1]{%
  \stepcounter{popsec}\setcounter{popsub}{0}%
  \par\vspace{16pt}\noindent
  \begin{minipage}{\linewidth}
  \tikz[baseline=-0.7ex]\node[draw=popink,line width=1.6pt,fill=poporange,rounded corners=4pt,minimum size=22pt,inner sep=0]{\popdisplay\fontsize{12}{12}\selectfont\color{popcream}\thepopsec};%
  \hspace{8pt}{\popdisplay\fontsize{15}{17}\selectfont\color{popink}\MakeUppercase{#1}}\par
  \vspace{4pt}\noindent{\color{poporange}\rule{36pt}{3.6pt}}
  \end{minipage}\par\nopagebreak\vspace{8pt}
}
\newcounter{popsub}
\newcommand{\courssub}[1]{%
  \stepcounter{popsub}%
  \par\vspace{9pt}\noindent{\popxbold\fontsize{11.5}{13}\selectfont\color{popink}{\color{poporange}\thepopsec.\thepopsub}\hspace{6pt}#1}\par\nopagebreak\vspace{4pt}
}

% ============================================================
% Boîtes pédagogiques — même recette : fond blanc, bordure épaisse colorée,
% ombre dure encre, badge plein. Titre optionnel [#1].
% ============================================================
\tcbset{popbase/.style={breakable,enhanced,colback=white,boxrule=2.3pt,arc=9pt,
  shadow={3pt}{-3pt}{0pt}{popink},left=14pt,right=14pt,top=11pt,bottom=11pt,before skip=11pt,after skip=15pt,
  fontupper=\popsemi\fontsize{10.6}{14.5}\selectfont\color{popink},
  fontlower=\popsemi\fontsize{10.6}{14.5}\selectfont\color{popink}}}
% Coche dessinée (le glyphe ✓ n'existe pas dans les polices du document)
\newcommand{\popcheck}[1]{\tikz[baseline=-0.5ex]\draw[#1,line width=1.6pt,line cap=round,line join=round](-0.13,0.0)--(-0.04,-0.09)--(0.13,0.1);}
\providecommand{\coche}{\popcheck{popgreen}}
\providecommand{\resultat}[1]{\par\smallskip\noindent\fcolorbox{popgreen}{popgreenL}{\parbox{\dimexpr\linewidth-2\fboxsep-2\fboxrule\relax}{\textbf{Résultat.} #1}}\par\smallskip}
\newcommand{\popboxhead}[3]{\poppill{#1}{#2}{white}\ifx\relax#3\relax\else\hspace{6pt}{\popxbold\fontsize{10.6}{12}\selectfont\color{popink}#3}\fi\par\vspace{7pt}}

\newtcolorbox{aretenir}[1][]{popbase,colframe=popgreen,before upper={\popboxhead{À retenir}{popgreen}{#1}}}
\newtcolorbox{exemplebox}[1][]{popbase,colframe=poporange,before upper={\popboxhead{Exemple}{poporange}{#1}}}
\newtcolorbox{definition}[1][]{popbase,colframe=popblue,before upper={\popboxhead{Définition}{popblue}{#1}}}
\newtcolorbox{propriete}[1][]{popbase,colframe=poppurple,before upper={\popboxhead{Propriété}{poppurple}{#1}}}
\newtcolorbox{methode}[1][]{popbase,colframe=popgold,before upper={\popboxhead{Méthode}{popgold}{#1}}}
\newtcolorbox{attention}[1][]{popbase,colframe=poppink,before upper={\popboxhead{Attention}{poppink}{#1}}}
\newtcolorbox{experience}[1][]{popbase,colframe=popblue,colback=popblueL,before upper={\popboxhead{Expérience}{popblue}{#1}}}
% « Le savais-tu ? » : carte violette pleine (contexte, culture, Cameroun)
\newtcolorbox{savaistu}[1][]{popbase,colback=poppurple,colframe=popink,
  fontupper=\popsemi\fontsize{10.4}{14.2}\selectfont\color{white},
  before upper={\poppill{Le savais-tu\,?}{white}{poppurple}\ifx\relax#1\relax\else\hspace{6pt}{\popxbold\color{white}#1}\fi\par\vspace{7pt}}}
% Exercice résolu : énoncé en haut, \tcblower puis la solution sur fond crème
\newcounter{popexo}\newcounter{popexr}
\newtcolorbox{exoresolu}[1][]{popbase,colframe=poporange,colbacklower=popcream,
  segmentation style={popink,line width=1.2pt,dashed},
  before upper={\stepcounter{popexr}\popboxhead{Exercice résolu \thepopexr}{poporange}{#1}},
  before lower={\poppill{Solution}{popink}{popcream}\par\vspace{6pt}}}
% Exercice (non résolu en place) — la correction est dans la partie Corrigés
\newtcolorbox{exercice}[1][]{popbase,colframe=popink,
  before upper={\stepcounter{popexo}\popboxhead{Exercice \thepopexo}{popink}{#1}}}
% Corrigé
\newtcolorbox{corrige}[1][]{popbase,colframe=popgreen,colback=popgreenL,
  before upper={\popboxhead{Corrigé}{popgreen}{#1}}}
% Objectifs / prérequis / bilan
\newtcolorbox{objectifs}{popbase,colframe=popink,colback=poporangeL,
  before upper={\popboxhead{Objectifs de la leçon}{popink}{}}}
\newtcolorbox{prerequis}{popbase,colframe=popink,colback=popblueL,
  before upper={\popboxhead{Prérequis}{popblue}{}}}
\newtcolorbox{bilan}{popbase,colframe=popink,colback=white,boxrule=2.8pt,
  before upper={\popboxhead{Fiche bilan}{poporange}{L'essentiel à connaître pour l'examen}}}
% Figure encadrée avec légende
\newtcolorbox{popfigure}[1][]{enhanced,breakable=false,colback=white,colframe=popline,boxrule=1.4pt,arc=8pt,
  left=10pt,right=10pt,top=10pt,bottom=8pt,before skip=10pt,after skip=12pt,halign=center,
  lower separated=false,halign lower=center,
  fontlower=\popsemi\fontsize{9}{11.5}\selectfont\color{popmuted},
  #1}
\newcommand{\legende}[1]{\par\vspace{6pt}{\popsemi\fontsize{9}{11.5}\selectfont\color{popmuted}#1}}

% ============================================================
% Signature OnBuch+
% ============================================================
\newcommand{\poplogoinline}[1]{{\poplogo\fontsize{#1}{#1}\selectfont\color{popink}OnBuch\color{poporange}+}}
\newcommand{\signatureonbuch}{%
  \begin{tcolorbox}[enhanced,colback=popink,colframe=popink,boxrule=2.2pt,arc=10pt,
      drop shadow=poporange,left=14pt,right=14pt,top=10pt,bottom=10pt,before skip=0pt,after skip=0pt]
    \tikz[baseline=-0.7ex]\node[circle,fill=popgreen,draw=popcream,line width=1.3pt,minimum size=22pt,inner sep=0]{\popcheck{white}};%
    \hspace{9pt}\begin{minipage}[c]{0.62\linewidth}
      {\popxbold\fontsize{12}{13}\selectfont\color{popcream}Signé\ \textbullet\ {\poplogo OnBuch\color{poporange}+}}\par\vspace{2pt}
      {\raggedright\popsemi\fontsize{8}{10.5}\selectfont\color{popline}\MakeUppercase{Équipe pédagogique OnBuch+}\\\MakeUppercase{Conforme au programme officiel MINESEC}\par}
    \end{minipage}\hfill
    {\poplogo\fontsize{22}{22}\selectfont\color{popcream}OnBuch\color{poporange}+}
  \end{tcolorbox}%
}

% ============================================================
% Page de garde
% #1 titre · #2 matière · #3 niveau · #4 module · #5 n° leçon · #6 durée ·
% #7 description / objectifs (texte ou itemize)
% ============================================================
\newcommand{\pagegarde}[7]{%
  \thispagestyle{empty}
  \newgeometry{a4paper,margin=1.7cm,top=1.6cm,bottom=1.4cm}%
  \begin{tikzpicture}[remember picture,overlay]
    \fill[poporange!18] ([xshift=-3.2cm,yshift=-3.6cm]current page.north east) circle (4.6cm);
    \fill[poppurple!14] ([xshift=2.4cm,yshift=7.6cm]current page.south west) circle (3.8cm);
  \end{tikzpicture}%
  {\poplogo\fontsize{30}{30}\selectfont\color{popink}OnBuch\color{poporange}+}\hfill
  \raisebox{0.6ex}{\poppill{Cours signé \textbullet\ Premium}{popink}{popcream}}\par
  \vspace{1pt}\popsquiggle{poppurple}\par
  \vspace{8pt}
  \begin{tcolorbox}[enhanced,colback=poporange,colframe=popink,boxrule=2.8pt,arc=13pt,
      drop shadow=popink,left=18pt,right=18pt,top=12pt,bottom=13pt,after skip=10pt]
    \poppill{#2 \textbullet\ #3}{popink}{popcream}\hspace{5pt}\poppill{#4}{white}{popink}\par\vspace{8pt}
    {\popdisplay\fontsize{14}{14}\selectfont\color{popink}LEÇON #5}\par\vspace{4pt}
    {\raggedright\hyphenpenalty=10000\exhyphenpenalty=10000\popdisplay\fontsize{25}{27}\selectfont\color{popcream}\MakeUppercase{#1}\par}
    \vspace{8pt}
    {\popxbold\fontsize{9.5}{9.5}\selectfont\color{popink}\MakeUppercase{Durée conseillée : #6}}
  \end{tcolorbox}
  \begin{tcolorbox}[enhanced,colback=white,colframe=popink,boxrule=2.2pt,arc=10pt,
      drop shadow=popink,left=16pt,right=16pt,top=10pt,bottom=10pt,after skip=8pt]
    \poppill{Dans cette leçon}{poppurple}{white}\par\vspace{5pt}
    \popsemi\fontsize{9.3}{12.6}\selectfont\color{popink}#7
  \end{tcolorbox}
  \noindent\begin{minipage}[t]{0.487\linewidth}
    \begin{tcolorbox}[enhanced,colback=popgreen,colframe=popink,boxrule=2.2pt,arc=10pt,
        drop shadow=popink,left=11pt,right=11pt,top=10pt,bottom=10pt,height=3.1cm,valign=center]
      \tcbox[colback=white,colframe=popink,boxrule=1.3pt,arc=5pt,boxsep=0pt,nobeforeafter,
        left=3pt,right=3pt,top=3pt,bottom=3pt,valign=center]{\qrcode[height=1.9cm]{https://play.google.com/store/apps/details?id=com.luvvix.onbuch&hl=fr}}%
      \hspace{8pt}\begin{minipage}[c]{0.5\linewidth}
        {\popdisplay\fontsize{11}{12.5}\selectfont\color{white}\MakeUppercase{Télécharge OnBuch}}\par\vspace{4pt}
        {\raggedright\popsemi\fontsize{8.4}{11}\selectfont\color{white}Gratuit \textbullet\ Google Play\\Tuteur IA, annales, concours, résultats\par}
      \end{minipage}
    \end{tcolorbox}
  \end{minipage}\hfill
  \begin{minipage}[t]{0.487\linewidth}
    \begin{tcolorbox}[enhanced,colback=poppurple,colframe=popink,boxrule=2.2pt,arc=10pt,
        drop shadow=popink,left=11pt,right=11pt,top=10pt,bottom=10pt,height=3.1cm,valign=center]
      \tikz[baseline=-0.7ex]\node[circle,fill=white,draw=popink,line width=1.3pt,minimum size=1.6cm,inner sep=0]{\popdisplay\fontsize{24}{24}\selectfont\color{poppurple}+};%
      \hspace{8pt}\begin{minipage}[c]{0.6\linewidth}
        {\popdisplay\fontsize{11}{12.5}\selectfont\color{white}\MakeUppercase{Passe à OnBuch+}}\par\vspace{4pt}
        {\raggedright\popsemi\fontsize{8.4}{11}\selectfont\color{white}Tous les cours signés, Tuteur IA illimité, fiches PDF — onglet OnBuch+ de l'app\par}
      \end{minipage}
    \end{tcolorbox}
  \end{minipage}
  \vspace{8pt}
  \popcontenu
  \vfill
  \signatureonbuch
  \restoregeometry
  \newpage
}

% Rangée « Ce cours contient » (page de garde)
\newcommand{\poptile}[3]{% #1 couleur #2 gros texte #3 légende
  \begin{tcolorbox}[enhanced,colback=#1,colframe=popink,boxrule=2pt,arc=9pt,shadow={2.5pt}{-2.5pt}{0pt}{popink},
      left=6pt,right=6pt,top=9pt,bottom=9pt,halign=center,valign=center,height=1.75cm,width=\linewidth,nobeforeafter]
    {\popdisplay\fontsize{12.5}{13}\selectfont\color{white}\MakeUppercase{#2}}\par\vspace{3pt}
    {\centering\popsemi\fontsize{7.8}{9.5}\selectfont\color{white}#3\par}
  \end{tcolorbox}}
\newcommand{\popcontenu}{%
  \par\noindent{\popxboldls\fontsize{7.5}{7.5}\selectfont\color{popmuted}CE COURS SIGNÉ CONTIENT}\par\vspace{7pt}
  \noindent\begin{minipage}[t]{0.235\linewidth}\poptile{popblue}{Cours}{complet, illustré, pas à pas}\end{minipage}\hfill
  \begin{minipage}[t]{0.235\linewidth}\poptile{poporange}{Méthodes}{et exercices résolus}\end{minipage}\hfill
  \begin{minipage}[t]{0.235\linewidth}\poptile{poppink}{Exercices}{type examen + corrigés}\end{minipage}\hfill
  \begin{minipage}[t]{0.235\linewidth}\poptile{popgreen}{Fiche bilan}{l'essentiel à retenir}\end{minipage}\par}

% Sommaire
\newtcolorbox{sommaire}{enhanced,colback=white,colframe=popink,boxrule=2.2pt,arc=10pt,
  drop shadow=popink,left=18pt,right=18pt,top=13pt,bottom=13pt,before skip=6pt,after skip=20pt,
  before upper={\poppill{Sommaire}{poppurple}{white}\par\vspace{9pt}\popsemi\fontsize{10.6}{14.5}\selectfont\color{popink}}}

% Signature de fin de cours — poussée tout en bas de la dernière page
\newcommand{\findecours}{%
  \par\vfill\nopagebreak
  \begin{center}\popsquiggle{poporange}\end{center}\vspace{4pt}
  \signatureonbuch
}
\AtBeginDocument{\def\llbracket{\mathopen{[\mkern-3.5mu[}}\def\rrbracket{\mathclose{]\mkern-3.5mu]}}} % intervalles d'entiers (glyphe ⟦ absent de la police)
% Glyphes absents de Fira Math : redéfinis à partir de glyphes disponibles
\AtBeginDocument{%
  \def\setminus{\mathbin{\backslash}}\let\smallsetminus\setminus
  \def\vdots{\mathord{\vbox{\baselineskip3.2pt\lineskiplimit0pt\kern4pt\hbox{.}\hbox{.}\hbox{.}}}}%
  \def\ddots{\mathinner{\mkern1mu\raise6.5pt\hbox{.}\mkern2mu\raise3.6pt\hbox{.}\mkern2mu\raise0.7pt\hbox{.}\mkern1mu}}%
  \def\triangle{\mathord{\scalebox{0.9}{$\symup{\Delta}$}}}%
  \def\nabla{\mathord{\raisebox{\depth}{\scalebox{1}[-1]{$\symup{\Delta}$}}}}%
  \def\checkmark{\popcheck{popdark}}%
  \def\star{\mathbin{\ast}}\def\bigstar{\mathord{\ast}}%
  \def\diamond{\mathbin{\scalebox{0.6}{$\Diamond$}}}%
  \def\bigcup{\mathop{\mathchoice{\raisebox{-0.3em}{\scalebox{1.7}{$\cup$}}}{\scalebox{1.25}{$\cup$}}{\cup}{\cup}}\displaylimits}%
  \def\bigcap{\mathop{\mathchoice{\raisebox{-0.3em}{\scalebox{1.7}{$\cap$}}}{\scalebox{1.25}{$\cap$}}{\cap}{\cap}}\displaylimits}%
  \def\lesseqqgtr{\mathrel{\lessgtr}}\def\bigoplus{\mathop{\oplus}\displaylimits}%
}
\providecommand{\ssi}{si et seulement si} % abréviation fréquente
\AtBeginDocument{\providecommand{\wideparen}{\overparen}} % arc de cercle

%% FIGFIX-BEGIN v5 — correctifs globaux des figures (docs/onbuch-generation/tools/figfix)
\usepackage{adjustbox}\usepackage{etoolbox}\tcbuselibrary{raster}
% toute figure TikZ/pgfplots plus large que la ligne est réduite à la largeur disponible
\BeforeBeginEnvironment{tikzpicture}{\begin{adjustbox}{max width=\linewidth}}
\AfterEndEnvironment{tikzpicture}{\end{adjustbox}}
% légendes pgfplots lisibles par-dessus les courbes
\pgfplotsset{every axis/.append style={legend style={fill=white,fill opacity=0.92,text opacity=1,draw=black!15,inner sep=3pt}}}
% étiquettes d'arêtes (syntaxe edge["..."]) avec fond blanc
\tikzset{every edge quotes/.append style={fill=white,inner sep=1.2pt}}
%% FIGFIX-END
\begin{document}

\pagegarde{Les pétroles}{PCT}{3ème}{Module 3 : Chimie appliquée}{N16}{2 séances (1 h chacune)}{Cette leçon présente le pétrole : sa définition, sa composition en hydrocarbures et les étapes de son raffinage. Elle décrit les principales utilisations des produits pétroliers et sensibilise à la lutte contre la pollution liée à leur exploitation et à leur consommation.
\vspace{4pt}
\begin{multicols}{2}\small
\begin{itemize}
\item À la fin de cette leçon, je sais définir le pétrole brut et citer ses principaux gisements au Cameroun
\item À la fin de cette leçon, je sais décrire la composition du pétrole en hydrocarbures
\item À la fin de cette leçon, je sais décrire les étapes du raffinage du pétrole, notamment la distillation fractionnée
\item À la fin de cette leçon, je sais interpréter le schéma d'une colonne de distillation fractionnée
\item À la fin de cette leçon, je sais citer les principaux produits pétroliers et leurs utilisations dans la vie courante
\item À la fin de cette leçon, je sais identifier les pollutions liées au pétrole et proposer des mesures de lutte
\end{itemize}\end{multicols}}

\begin{sommaire}
\begin{multicols}{2}
\begin{enumerate}
\item Le pétrole brut : définition, origine et gisements
\item Composition du pétrole : les hydrocarbures
\item Le raffinage du pétrole : la distillation fractionnée
\item Les produits pétroliers et leurs utilisations
\item Pétrole et pollution : impacts et lutte
\item Activité d'intégration
\item Méthodes et exercices résolus
\item Exercices
\item Corrigés des exercices
\item Fiche bilan
\end{enumerate}
\end{multicols}
\end{sommaire}

\begin{objectifs}
\begin{itemize}
\item À la fin de cette leçon, je sais définir le pétrole brut et citer ses principaux gisements au Cameroun
\item À la fin de cette leçon, je sais décrire la composition du pétrole en hydrocarbures
\item À la fin de cette leçon, je sais décrire les étapes du raffinage du pétrole, notamment la distillation fractionnée
\item À la fin de cette leçon, je sais interpréter le schéma d'une colonne de distillation fractionnée
\item À la fin de cette leçon, je sais citer les principaux produits pétroliers et leurs utilisations dans la vie courante
\item À la fin de cette leçon, je sais identifier les pollutions liées au pétrole et proposer des mesures de lutte
\item À la fin de cette leçon, je sais appliquer les consignes de sécurité lors de la manipulation et du stockage des produits pétroliers
\item À la fin de cette leçon, je sais exploiter un tableau ou un graphique de données sur la production ou les fractions du pétrole
\end{itemize}
\end{objectifs}

\begin{prerequis}
\begin{itemize}
\item Les hydrocarbures et la formule générale des alcanes (leçon sur les hydrocarbures, 3ème)
\item La combustion complète et incomplète des hydrocarbures
\item Les changements d'état : vaporisation, liquéfaction, température d'ébullition (classe de 4ème)
\item La notion de mélange et les techniques de séparation (distillation simple vue en 4ème)
\item Les combustions et leurs effets sur l'environnement (gaz à effet de serre, notions vues en 4ème)
\end{itemize}
\end{prerequis}

\newpage

\coursec{Le pétrole brut : définition, origine et gisements}

\textbf{Situation d'accroche.} Il est 6 h 30 à Douala. Déjà, des milliers de motos-taxis vrombissent sur le boulevard de la Liberté, les bus de ligne chargent leurs passagers, et dans les quartiers, des ménagères allument leurs cuisinières au gaz butane pour préparer le repas. Essence, gasoil, gaz butane : tous ces produits, si différents en apparence, proviennent d'une seule et même matière noire, épaisse et visqueuse, extraite du sous-sol camerounais — notamment au large de Kribi et dans le bassin du Rio del Rey, près de Limbé. Cette matière, c'est le \cle{pétrole brut}. Pourtant, personne ne verse du pétrole brut dans le réservoir d'une moto ni dans une bonbonne de gaz : tel quel, il est inutilisable. Alors, d'où vient cette matière noire ? Comment s'est-elle formée et où se cache-t-elle dans le sous-sol ? Et comment la transforme-t-on en produits utilisables ? C'est à ces questions que cette leçon va répondre, pas à pas.

\courssub{Qu'est-ce que le pétrole brut ?}

Observe une bouteille d'huile de vidange usagée sortant d'un moteur de moto-taxi : un liquide noir, épais, qui coule lentement et colle aux doigts. Le pétrole brut ressemble à cela, mais en plus complexe. Ce n'est pas une substance pure : c'est un \cle{mélange}.

\begin{definition}[Pétrole brut]
Le \cle{pétrole brut} est un mélange complexe d'\cle{hydrocarbures} liquides, dans lesquels sont dissous des hydrocarbures gazeux et solides. C'est un liquide d'aspect \textbf{noir} (parfois brun ou verdâtre), \textbf{visqueux}, \textbf{inflammable} et \textbf{insoluble dans l'eau}, sur laquelle il flotte car il est moins dense qu'elle.
\end{definition}

Retenons trois idées essentielles :

\begin{itemize}
\item Le mot \emph{pétrole} vient du latin \emph{petra} (pierre) et \emph{oleum} (huile) : littéralement, « huile de pierre », car il suinte parfois naturellement à la surface des roches.
\item Le mot \emph{brut} signifie « non traité, non transformé » : c'est le pétrole tel qu'il sort du sol, avant toute transformation.
\item Un \cle{hydrocarbure} est un composé chimique formé uniquement d'atomes de \cle{carbone} (C) et d'\cle{hydrogène} (H). Nous y reviendrons en détail dans la section 2.
\end{itemize}

\begin{attention}[Erreur fréquente : pétrole brut et essence]
Ne confonds jamais le \textbf{pétrole brut} et l'\textbf{essence}. Le pétrole brut est la matière première extraite du sous-sol ; l'essence n'est qu'\textbf{un des produits} obtenus après transformation du brut à la raffinerie. Dire « ma moto roule au pétrole brut » est faux : elle roule à l'essence, qui provient du raffinage du pétrole brut. De même, le gasoil, le kérosène et le gaz butane sont des \emph{produits pétroliers}, pas du pétrole brut.
\end{attention}

\begin{exemplebox}[Observation de la vie courante]
À la station-service, on vend de l'essence (liquide clair, jaunâtre, très fluide) et du gasoil (liquide jaune ambré). Aucun de ces carburants ne ressemble au liquide noir et épais qui sort d'un puits de pétrole. Cette différence d'aspect est une première preuve que le brut a subi des transformations avant d'arriver à la pompe.
\end{exemplebox}

\courssub{L'origine du pétrole : une formation très lente}

D'où vient cette huile cachée dans les roches ? Pour le comprendre, menons une petite démarche d'investigation à partir de documents.

\begin{experience}[Investigation : que contiennent les roches pétrolifères ?]
\textbf{Document 1.} Les géologues qui étudient les roches d'où provient le pétrole (appelées \cle{roches mères}) y trouvent de nombreux fossiles de microorganismes marins : restes de plancton, d'algues microscopiques, de petits coquillages.

\textbf{Document 2.} Ces roches se sont formées au fond d'anciennes mers, recouvertes ensuite par d'épaisses couches de sédiments (sable, argile) sur plusieurs kilomètres d'épaisseur.

\textbf{Document 3.} À de grandes profondeurs, la température et la pression sont élevées.

\textbf{Hypothèse.} Le pétrole proviendrait de la transformation de matières organiques enfouies.

\textbf{Interprétation.} Les matières organiques mortes (plancton, végétaux) se déposent au fond des mers avec les sédiments. Enfouies de plus en plus profondément au fil du temps, elles subissent une pression forte et une température élevée, en l'absence d'air. Leur décomposition lente produit des hydrocarbures : le pétrole et le gaz naturel.

\textbf{Conclusion.} Le modèle actuellement admis est celui de l'origine \cle{organique} du pétrole : il résulte de la décomposition de matières organiques enfouies pendant des millions d'années, sous l'effet de la pression et de la température.
\end{experience}

\begin{aretenir}[Formation du pétrole : les étapes]
\begin{enumerate}
\item Des matières organiques (plancton, algues, végétaux) meurent et se déposent au fond des mers ou des lagunes, mélangées aux sédiments.
\item Les couches de sédiments s'accumulent et enfoncent cette matière organique, à l'abri de l'air.
\item Sous l'effet de la \textbf{pression} et de la \textbf{température}, pendant des \textbf{millions d'années}, la matière organique se transforme lentement en hydrocarbures liquides (pétrole) et gazeux (gaz naturel) dans la \cle{roche mère}.
\item Le pétrole, plus léger que l'eau contenue dans les roches, migre lentement vers le haut jusqu'à être bloqué par une roche imperméable : il s'accumule alors dans un \cle{gisement}.
\end{enumerate}
\end{aretenir}

\begin{exemplebox}[Pourquoi dit-on que le pétrole est une énergie non renouvelable ?]
La formation du pétrole demande des \textbf{millions d'années} (on estime généralement de 10 à 400 millions d'années). Or l'humanité consomme les gisements en quelques dizaines d'années. Le stock ne peut donc pas se reconstituer à l'échelle d'une vie humaine : on dit que le pétrole est une \cle{énergie fossile} et \cle{non renouvelable}. Chaque litre d'essence brûlé aujourd'hui est un héritage de la préhistoire qui ne reviendra jamais.
\end{exemplebox}

\courssub{Gisement, piège pétrolier et extraction}

Le pétrole ne forme pas de « lacs » souterrains : il est emprisonné dans les \textbf{pores} (les petits trous) d'une roche, comme l'eau dans une éponge. Pour qu'un gisement exploitable se forme, il faut un arrangement particulier des couches de roches, appelé \cle{piège pétrolier}.

\begin{definition}[Gisement et piège pétrolier]
Un \cle{gisement de pétrole} est une accumulation de pétrole dans les pores d'une roche perméable appelée \cle{roche réservoir}. Le gisement n'existe que si un \cle{piège} le retient : une \cle{roche couverture} imperméable (souvent de l'argile), bombée ou inclinée, qui empêche le pétrole de remonter vers la surface.
\end{definition}

Dans un piège typique, on trouve de haut en bas :
\begin{itemize}
\item le \textbf{gaz naturel}, le plus léger, tout en haut du piège ;
\item le \textbf{pétrole brut}, au milieu ;
\item l'\textbf{eau salée} (saumure), en bas, car elle est la plus dense.
\end{itemize}

\begin{popfigure}
\begin{center}
\begin{tikzpicture}[scale=1.0]
% couches de fond
\fill[popgoldL] (-6,-3.4) rectangle (6,-2.2);
\fill[popmuted!40] (-6,-2.2) rectangle (6,-1.2);
% roche couverture bombee
\fill[popgreenL] plot[smooth] coordinates {(-6,-1.2) (-3,-0.4) (0,0.1) (3,-0.4) (6,-1.2)} -- (6,-0.2) -- (-6,-0.2) -- cycle;
\fill[popgreenL] (-6,-0.2) rectangle (6,0.4);
% roche reservoir (piege)
\fill[poporangeL] plot[smooth] coordinates {(-3.6,-1.2) (-2,-0.75) (0,-0.35) (2,-0.75) (3.6,-1.2)} -- (3.6,-2.2) -- (-3.6,-2.2) -- cycle;
% gaz
\fill[poppinkL] plot[smooth] coordinates {(-1.7,-0.78) (0,-0.42) (1.7,-0.78)} -- (1.7,-1.25) -- (-1.7,-1.25) -- cycle;
% petrole
\fill[popdark!85] plot[smooth] coordinates {(-3.2,-1.25) (0,-1.25) (3.2,-1.25)} -- (3.2,-1.75) -- (-3.2,-1.75) -- cycle;
% eau
\fill[popblueL] (-3.6,-1.75) rectangle (3.6,-2.2);
% surface et puits
\fill[popcream] (-6,0.4) rectangle (6,1.0);
\draw[popdark,thick] (-6,0.4) -- (6,0.4);
\draw[line width=3pt,popdark] (0.6,1.0) -- (0.6,-0.5);
\draw[popdark,fill=popdark] (0.6,1.0) circle (0.06);
% derrick simple
\draw[popdark,thick] (0.2,1.0) -- (0.6,2.2) -- (1.0,1.0);
\draw[popdark,thick] (0.32,1.5) -- (0.88,1.5);
% legendes flechees
\node[font=\small,anchor=west] at (3.9,0.1) {\textbf{1}};
\draw[->,popdark] (3.85,0.1) -- (2.6,-0.1);
\node[font=\small,anchor=west] at (3.9,-0.95) {\textbf{2}};
\draw[->,popdark] (3.85,-0.95) -- (1.2,-0.85);
\node[font=\small,anchor=west] at (3.9,-1.5) {\textbf{3}};
\draw[->,popdark] (3.85,-1.5) -- (1.5,-1.5);
\node[font=\small,anchor=west] at (3.9,-1.98) {\textbf{4}};
\draw[->,popdark] (3.85,-1.98) -- (1.5,-1.95);
\node[font=\small,anchor=west] at (3.9,-2.8) {\textbf{5}};
\draw[->,popdark] (3.85,-2.8) -- (1.0,-2.8);
\node[font=\small,anchor=east] at (-1.2,1.7) {\textbf{6}};
\draw[->,popdark] (-1.15,1.65) -- (0.5,1.5);
\end{tikzpicture}
\end{center}
\legende{Coupe schématique d'un gisement pétrolier. \textbf{1} : roche couverture imperméable (argile), bombée, qui forme le piège. \textbf{2} : gaz naturel (le plus léger). \textbf{3} : pétrole brut emprisonné dans les pores de la roche réservoir. \textbf{4} : eau salée. \textbf{5} : roche mère, où le pétrole s'est formé. \textbf{6} : puits de forage avec son derrick, qui traverse les couches pour atteindre le gisement.}
\end{popfigure}

\begin{definition}[Forage]
Le \cle{forage} est l'opération qui consiste à creuser un puits étroit et profond (parfois plus de \SI{3}{km}) à l'aide d'un trépan rotatif, afin d'atteindre la roche réservoir et d'extraire le pétrole brut. En mer, le forage est réalisé depuis des \textbf{plates-formes pétrolières} : on parle d'exploitation \emph{offshore}.
\end{definition}

Lorsque le puits atteint le gisement, la pression naturelle fait parfois jaillir le pétrole ; ensuite, on utilise des pompes (les fameux « chevalets de pompage » que l'on voit dans certaines régions) ou on injecte de l'eau pour pousser le brut vers la surface.

\courssub{Les gisements du Cameroun et la SONARA}

Le Cameroun est un pays producteur de pétrole depuis 1977. Ses principaux gisements se situent dans trois bassins sédimentaires :

\begin{itemize}
\item le \textbf{bassin du Rio del Rey}, au large de la région du Sud-Ouest (zone de Limbé et Bakassi) : c'est le bassin historique, le plus ancien exploité ;
\item le \textbf{bassin de Douala/Kribi-Campo}, au large des côtes entre Douala, Kribi et Campo : il fournit aujourd'hui l'essentiel de la production, avec des gisements offshore ;
\item le \textbf{bassin de Logone-Birni}, dans l'Extrême-Nord, à la frontière du Tchad : bassin terrestre dont l'exploitation est plus récente et plus modeste.
\end{itemize}

\begin{popfigure}
\begin{center}
\begin{tikzpicture}[scale=0.9]
% contour simplifie du Cameroun
\draw[popdark,thick,fill=popcream] plot[smooth cycle] coordinates {
(1.2,0.2) (0.4,1.2) (0.6,2.6) (1.0,3.6) (0.8,4.6) (1.6,5.4)
(1.4,6.4) (2.2,7.2) (3.4,7.8) (4.6,7.4) (5.6,6.6) (6.2,5.6)
(6.6,4.4) (6.0,3.4) (6.4,2.4) (5.6,1.4) (4.4,0.8) (3.0,0.4) (2.0,-0.2)};
% mer
\fill[popblueL] (-1.6,-0.6) rectangle (0.6,3.4);
\node[font=\small\itshape,popblue] at (-0.5,2.6) {Golfe};
\node[font=\small\itshape,popblue] at (-0.5,2.2) {de Guinée};
% bassins
\fill[poporange,opacity=0.55] (0.15,1.5) circle (0.55);
\fill[popgold,opacity=0.6] (0.5,0.35) circle (0.5);
\fill[popgreen,opacity=0.5] (3.4,6.6) circle (0.7);
% raffinerie Limbe
\node[star,star points=5,fill=popink,inner sep=1.5pt] at (0.55,1.15) {};
% villes
\fill[popdark] (1.5,1.1) circle (0.07); \node[font=\small,anchor=west] at (1.6,1.1) {Douala};
\fill[popdark] (1.35,0.15) circle (0.07); \node[font=\small,anchor=west] at (1.45,0.1) {Kribi};
\node[font=\small,anchor=west] at (0.9,1.45) {Limbé};
% legendes
\node[font=\small,anchor=west,poporange!80!black] at (2.2,2.6) {Bassin du Rio del Rey};
\draw[->,poporange!80!black] (2.2,2.5) -- (0.5,1.7);
\node[font=\small,anchor=west,popgold!70!black] at (2.2,2.0) {Bassin Douala/Kribi-Campo};
\draw[->,popgold!70!black] (2.2,1.9) -- (0.85,0.45);
\node[font=\small,anchor=west,popgreen!60!black] at (4.3,6.9) {Bassin de Logone-Birni};
\draw[->,popgreen!60!black] (4.3,6.8) -- (3.8,6.55);
\node[font=\small,anchor=west] at (2.2,3.2) {\textcolor{popink}{$\bigstar$} Raffinerie SONARA (Limbé)};
\end{tikzpicture}
\end{center}
\legende{Carte simplifiée du Cameroun localisant les trois bassins pétroliers (Rio del Rey, Douala/Kribi-Campo, Logone-Birni) et la raffinerie de la SONARA à Limbé. La plupart des gisements sont offshore, c'est-à-dire en mer, dans le golfe de Guinée.}
\end{popfigure}

\begin{savaistu}[La SONARA, la raffinerie de Limbé]
La \cle{SONARA} (Société Nationale de Raffinage), installée à Limbé au pied du mont Cameroun, est la \textbf{seule raffinerie du Cameroun}. Inaugurée en 1981, elle peut traiter environ 2 millions de tonnes de brut par an. Fait étonnant : une grande partie du pétrole brut camerounais est \textbf{exportée directement}, car il est de très bonne qualité et très recherché ; la SONARA raffine donc aussi du brut \textbf{importé} (notamment du Nigeria), mieux adapté à ses installations. Le Cameroun exporte du brut et importe une partie de ses carburants : voilà pourquoi les prix à la pompe dépendent du marché mondial.
\end{savaistu}

\begin{exemplebox}[Quelques données chiffrées]
\begin{itemize}
\item Le pétrole représente environ \textbf{40 \%} des recettes d'exportation du Cameroun : c'est l'un des premiers produits vendus par le pays à l'étranger.
\item Les hydrocarbures (pétrole et gaz) fournissent une part importante de l'énergie consommée au Cameroun, notamment pour le transport et la production d'électricité (centrales thermiques, groupes électrogènes).
\item La production camerounaise est de l'ordre de quelques dizaines de millions de barils par an (1 baril $\approx$ 159 litres).
\end{itemize}
Ces chiffres montrent que le pétrole est une richesse stratégique pour le pays — mais aussi une richesse qui s'épuise.
\end{exemplebox}

\courssub{Le pétrole brut n'est pas utilisable directement}

Verserais-tu dans le réservoir d'une moto un liquide noir, épais, contenant à la fois des gaz, des liquides légers et des matières solides ? Évidemment non : le moteur serait encrassé et détruit en quelques minutes.

\begin{aretenir}[Pourquoi faut-il raffiner le pétrole brut ?]
Le pétrole brut est un \textbf{mélange} d'hydrocarbures très différents : gaz légers, liquides fluides, liquides épais, solides (paraffines, bitume), ainsi que des impuretés (soufre, sels, eau). Aucun appareil — moteur, cuisinière, lampe — ne peut utiliser ce mélange tel quel. Il faut donc \textbf{séparer} les constituants du brut et les \textbf{traiter} : c'est le rôle du \cle{raffinage}, réalisé dans une raffinerie comme la SONARA. La technique principale est la \cle{distillation fractionnée}, que nous étudierons dans la section 3.
\end{aretenir}

\begin{exoresolu}[Un raisonnement sur le gisement]
\textbf{Énoncé.} Lors d'un forage dans le bassin du Rio del Rey, le puits traverse d'abord une couche d'argile, puis atteint une roche poreuse contenant du gaz, du pétrole et de l'eau salée. 1) Comment appelle-t-on la couche d'argile ? Quel est son rôle ? 2) Dans quel ordre (du haut vers le bas) se trouvent le gaz, le pétrole et l'eau dans la roche poreuse ? Justifie. 3) Pourquoi le pétrole extrait ne peut-il pas être versé directement dans le réservoir d'un bus ?
\tcblower
\textbf{Solution détaillée.}

1) La couche d'argile est la \textbf{roche couverture}. L'argile est \textbf{imperméable} : elle bloque la remontée du pétrole et du gaz vers la surface et forme ainsi le \textbf{piège pétrolier}, qui a permis l'accumulation du gisement.

2) De haut en bas, on trouve : le \textbf{gaz naturel}, puis le \textbf{pétrole brut}, puis l'\textbf{eau salée}. Justification : les fluides se rangent selon leur densité ; le gaz est le moins dense, l'eau salée est la plus dense, le pétrole occupe une position intermédiaire (il flotte sur l'eau).

3) Le pétrole extrait est du \textbf{pétrole brut} : un mélange complexe d'hydrocarbures liquides, gazeux et solides, avec des impuretés. Il est trop épais, trop encrassant et inadapté à un moteur. Il doit d'abord être \textbf{raffiné} (distillation fractionnée et traitements) pour obtenir le gasoil qui alimentera le bus.
\end{exoresolu}

\begin{attention}[Pièges de rédaction au BEPC]
\begin{itemize}
\item Ne dis pas que le pétrole forme des « rivières » ou des « lacs » souterrains : il est contenu dans les \textbf{pores} d'une roche réservoir.
\item La formation du pétrole n'est pas un fait observé directement : c'est un \textbf{modèle admis} à partir d'indices (fossiles, roches sédimentaires). Rédige : « on admet que le pétrole provient de la décomposition de matières organiques enfouies ».
\item Distingue bien les trois roches : \textbf{roche mère} (lieu de formation), \textbf{roche réservoir} (lieu de stockage), \textbf{roche couverture} (couvercle imperméable du piège).
\end{itemize}
\end{attention}

\begin{exercice}[Pour t'entraîner]
1) Donne la définition du pétrole brut et cite trois de ses propriétés physiques.

2) Explique en quatre étapes comment le pétrole s'est formé selon le modèle actuellement admis.

3) Pourquoi dit-on que le pétrole est une énergie non renouvelable ?

4) Cite les trois bassins pétroliers du Cameroun et le nom de la ville où se trouve la raffinerie nationale.

5) Un élève affirme : « La SONARA extrait le pétrole au large de Kribi ». Corrige cette affirmation.
\end{exercice}

\begin{corrige}[Exercice]
1) Le pétrole brut est un mélange complexe d'hydrocarbures liquides, gazeux et solides dissous. Propriétés : il est noir, visqueux, inflammable, insoluble dans l'eau et moins dense qu'elle (trois propriétés suffisent).

2) Dépôt de matières organiques mortes (plancton, végétaux) au fond des mers ; enfouissement sous les sédiments ; transformation sous l'effet de la pression et de la température pendant des millions d'années dans la roche mère ; migration et accumulation dans un piège (roche réservoir coiffée d'une roche couverture).

3) Parce que sa formation demande des millions d'années alors que sa consommation est très rapide : le stock ne se reconstitue pas à l'échelle humaine.

4) Bassin du Rio del Rey, bassin de Douala/Kribi-Campo, bassin de Logone-Birni ; la raffinerie (SONARA) se trouve à Limbé.

5) La SONARA ne fait pas d'extraction : c'est une \textbf{raffinerie} installée à Limbé, qui transforme le pétrole brut en produits finis. L'extraction au large de Kribi est réalisée par des compagnies d'exploitation (forages offshore).
\end{corrige}

En résumé, le pétrole brut est un mélange d'hydrocarbures né de la lente transformation de matières organiques, piégé dans les roches du sous-sol et extrait par forage. Mais ce trésor noir reste inutilisable en l'état. Dans la section suivante, nous examinerons de près sa composition : que sont exactement ces fameux hydrocarbures ?

\coursec{Composition du pétrole : les hydrocarbures}

Dans la section précédente, nous avons découvert que le pétrole brut est un liquide noir et visqueux, formé il y a des millions d'années à partir d'organismes marins, et piégé dans des roches réservoirs, comme celles du bassin du Rio del Rey au Cameroun. Mais de quoi ce liquide est-il réellement constitué ? Pourquoi peut-on en tirer à la fois du gaz butane, de l'essence pour moto-taxi et du gasoil pour camion ? Pour le comprendre, il faut pénétrer dans la composition chimique du pétrole : c'est l'objet de cette section.

\courssub{Le pétrole : un mélange d'hydrocarbures}

\textbf{Une observation pour commencer.} Verse dans un verre un peu d'huile de cuisine et de l'eau : les deux liquides ne se mélangent pas, mais chacun d'eux est homogène. Maintenant, observe une flaque d'essence sur le sol : elle s'évapore en quelques minutes, alors qu'une flaque de gasoil disparaît beaucoup plus lentement. Pourtant, essence et gasoil proviennent du même pétrole brut ! Cela signifie que le pétrole contient des substances différentes, aux propriétés différentes. Le pétrole n'est donc pas une substance unique : c'est un \cle{mélange}.

\begin{definition}[Hydrocarbure]
Un \cle{hydrocarbure} est un composé chimique dont la molécule est formée \textbf{uniquement} d'atomes de \textbf{carbone} (C) et d'atomes d'\textbf{hydrogène} (H).
\end{definition}

Le pétrole brut est un \textbf{mélange complexe} de très nombreux hydrocarbures, dont les molécules contiennent de 1 à plus de 40 atomes de carbone. C'est cette diversité de molécules qui explique la diversité des produits que l'on en tire.

\begin{attention}[Le pétrole n'est pas un corps pur]
Erreur très fréquente au BEPC : écrire « le pétrole est un corps pur composé » ou « le pétrole a une formule chimique ». C'est \textbf{faux} ! Un corps pur possède une formule chimique unique (comme l'eau, \ce{H2O}) et une température d'ébullition fixe. Le pétrole est un \textbf{mélange} d'hydrocarbures : il n'a \textbf{pas de formule chimique} et il ne bout pas à une température précise, mais sur une large plage de températures. Retiens : pétrole = mélange.
\end{attention}

\courssub{Les alcanes : la famille principale du pétrole}

\textbf{Rappel de la leçon sur les alcanes.} Les \cle{alcanes} sont des hydrocarbures dont les atomes de carbone sont reliés par des liaisons simples. Ils obéissent tous à la même \textbf{formule générale} :
\[ \ce{C_{n}H_{2n+2}} \]
où $n$ est le nombre d'atomes de carbone de la molécule ($n \geq 1$). Pour trouver le nombre d'atomes d'hydrogène, on multiplie $n$ par 2 et on ajoute 2.

\begin{exemplebox}[Quelques alcanes présents dans le pétrole]
\begin{itemize}
\item Le \textbf{méthane} \ce{CH4} ($n = 1$) : un seul atome de carbone, $2 \times 1 + 2 = 4$ atomes d'hydrogène. C'est le principal constituant du gaz naturel.
\item Le \textbf{butane} \ce{C4H10} ($n = 4$) : $2 \times 4 + 2 = 10$ atomes d'hydrogène. C'est le gaz des bonbonnes de cuisine.
\item L'\textbf{octane} \ce{C8H18} ($n = 8$) : $2 \times 8 + 2 = 18$ atomes d'hydrogène. C'est un constituant typique de l'essence.
\end{itemize}
\end{exemplebox}

\begin{methode}[Vérifier qu'une molécule est un alcane]
\begin{enumerate}
\item Compte le nombre $n$ d'atomes de carbone dans la formule.
\item Calcule $2n + 2$.
\item Compare avec le nombre d'atomes d'hydrogène de la formule : s'ils sont égaux, c'est un alcane ; sinon, ce n'en est pas un.
\end{enumerate}
Exemple : \ce{C6H14} donne $2 \times 6 + 2 = 14$ : c'est bien un alcane (l'hexane). En revanche \ce{C6H12} donne $2 \times 6 + 2 = 14 \neq 12$ : ce n'est pas un alcane.
\end{methode}

\courssub{Les autres constituants du pétrole}

Le pétrole brut ne contient pas que des alcanes. On y trouve aussi, en plus petites proportions :

\begin{itemize}
\item des \cle{cycloalcanes} : des hydrocarbures dont les atomes de carbone forment un \textbf{cycle} (un anneau), comme le cyclohexane \ce{C6H12} ;
\item des \cle{hydrocarbures aromatiques} : des hydrocarbures contenant le noyau benzénique (cycle de 6 atomes de carbone particulier), comme le benzène \ce{C6H6} ;
\item des \textbf{impuretés} : des composés contenant du \textbf{soufre} (S) ou de l'\textbf{azote} (N). Ces impuretés sont gênantes : à la combustion, elles produisent des gaz acides polluants (dioxyde de soufre \ce{SO2}, oxydes d'azote). C'est pourquoi les raffineries, comme la SONARA à Limbé, cherchent à les éliminer (désulfuration).
\end{itemize}

\begin{aretenir}[Composition du pétrole brut]
Le pétrole brut est un \textbf{mélange} constitué essentiellement d'\textbf{hydrocarbures} : des alcanes ($\ce{C_nH_{2n+2}}$), des cycloalcanes et des hydrocarbures aromatiques. Il contient aussi des impuretés (composés soufrés et azotés).
\end{aretenir}

\courssub{Propriété clé : la température d'ébullition augmente avec le nombre d'atomes de carbone}

\textbf{Question d'investigation.} Le gaz butane s'échappe de la bonbonne dès qu'on ouvre le robinet, alors que l'huile de vidange d'une moto reste liquide et épaisse même par forte chaleur. Toutes deux viennent du pétrole. D'où vient cette différence de comportement ? Formulons une hypothèse : \emph{plus la molécule d'hydrocarbure est grosse (plus elle contient d'atomes de carbone), plus il faut la chauffer pour la faire bouillir.}

Pour vérifier cette hypothèse, exploitons les données expérimentales mesurées au laboratoire pour quelques alcanes.

\begin{center}
\begin{tabularx}{\linewidth}{|l|c|c|X|}
\hline
\rowcolor{popblue!40} \textbf{Alcane} & \textbf{Formule} & \textbf{Nombre $n$ de C} & \textbf{Température d'ébullition} \\
\hline
Méthane & \ce{CH4} & 1 & \SI{-162}{\celsius} \\
\hline
Butane & \ce{C4H10} & 4 & \SI{-0,5}{\celsius} \\
\hline
Octane & \ce{C8H18} & 8 & \SI{126}{\celsius} \\
\hline
Hexadécane & \ce{C16H34} & 16 & \SI{287}{\celsius} \\
\hline
\end{tabularx}
\end{center}

\textbf{Observation.} En lisant le tableau de haut en bas, on constate que lorsque le nombre d'atomes de carbone augmente (1, 4, 8, 16), la température d'ébullition augmente aussi ($-162$, $-0,5$, $126$, \SI{287}{\celsius}). Notre hypothèse semble confirmée. Pour rendre cette évolution plus visible, traçons la courbe.

\begin{methode}[Tracer une courbe d'évolution à partir d'un tableau de données]
\begin{enumerate}
\item \textbf{Choisir les axes} : la grandeur que l'on fait varier (ici le nombre $n$ d'atomes de carbone) se porte en \textbf{abscisse} (axe horizontal) ; la grandeur mesurée (ici la température d'ébullition) se porte en \textbf{ordonnée} (axe vertical).
\item \textbf{Choisir une échelle} adaptée : par exemple 1 cm pour 2 atomes de carbone en abscisse et 1 cm pour \SI{50}{\celsius} en ordonnée, pour que la courbe occupe bien la feuille.
\item \textbf{Placer les points} : chaque ligne du tableau donne un point de coordonnées $(n \, ; \theta_{\text{éb}})$.
\item \textbf{Relier les points} par une courbe régulière (pas des segments de droite si le phénomène est continu).
\item \textbf{Légender} : écrire la grandeur et l'unité sur chaque axe, et donner un titre à la courbe.
\end{enumerate}
\end{methode}

\begin{popfigure}
\begin{center}
\begin{tikzpicture}
\begin{axis}[
width=0.85\linewidth,
height=7cm,
xlabel={Nombre d'atomes de carbone $n$},
ylabel={Température d'ébullition (\si{\celsius})},
xmin=0, xmax=18,
ymin=-200, ymax=350,
xtick={0,2,4,6,8,10,12,14,16,18},
ytick={-200,-100,0,100,200,300},
grid=both,
grid style={popline!40},
legend pos=south east,
legend style={font=\small}
]
\addplot[popcurve,smooth,tension=0.6] coordinates {(1,-162) (4,-0.5) (8,126) (16,287)};
\addlegendentry{Courbe d'évolution}
\addplot[only marks,mark=*,mark size=2.5pt,poporange] coordinates {(1,-162) (4,-0.5) (8,126) (16,287)};
\addlegendentry{Points expérimentaux}
\node[font=\small,pin={90:méthane}] at (axis cs:1,-162) {};
\node[font=\small,pin={270:butane}] at (axis cs:4,-0.5) {};
\node[font=\small,pin={180:octane}] at (axis cs:8,126) {};
\node[font=\small,pin={90:hexadécane}] at (axis cs:16,287) {};
\end{axis}
\end{tikzpicture}
\end{center}
\legende{Température d'ébullition des alcanes en fonction du nombre d'atomes de carbone $n$. La courbe est croissante : plus la chaîne carbonée est longue, plus la température d'ébullition est élevée.}
\end{popfigure}

\textbf{Interprétation.} La courbe obtenue est \textbf{croissante} : elle monte régulièrement de gauche à droite. Les petites molécules (peu de carbone) bouillent à basse température : ce sont des gaz à température ordinaire (méthane, butane). Les grosses molécules (beaucoup de carbone) bouillent à haute température : ce sont des liquides épais, voire des solides.

\begin{propriete}[Température d'ébullition des alcanes]
La \textbf{température d'ébullition} d'un alcane \textbf{augmente} lorsque le \textbf{nombre d'atomes de carbone} de sa molécule augmente (on dit aussi : lorsque la longueur de la chaîne carbonée augmente). Cette propriété, établie expérimentalement, se généralise aux autres hydrocarbures du pétrole.
\end{propriete}

\begin{exoresolu}[Exploiter la propriété]
Énoncé. On considère trois hydrocarbures extraits du pétrole : l'éthane \ce{C2H6}, le pentane \ce{C5H12} et le décane \ce{C10H22}. 1) Vérifie que ce sont des alcanes. 2) Sans connaître les valeurs exactes, classe ces trois corps par température d'ébullition croissante. Justifie. 3) Lequel est probablement un gaz à la température ordinaire de \SI{25}{\celsius} ?
\tcblower
Solution détaillée.
\begin{enumerate}
\item Appliquons la formule générale $\ce{C_nH_{2n+2}}$ :
\begin{itemize}
\item éthane : $n = 2$, donc $2n + 2 = 6$ : \ce{C2H6} est bien un alcane ;
\item pentane : $n = 5$, donc $2n + 2 = 12$ : \ce{C5H12} est bien un alcane ;
\item décane : $n = 10$, donc $2n + 2 = 22$ : \ce{C10H22} est bien un alcane.
\end{itemize}
\item La température d'ébullition d'un alcane croît avec le nombre d'atomes de carbone. Or $2 < 5 < 10$. Donc :
\[ \theta_{\text{éb}}(\text{éthane}) < \theta_{\text{éb}}(\text{pentane}) < \theta_{\text{éb}}(\text{décane}) \]
\item L'éthane, qui a la plus petite molécule, a la température d'ébullition la plus basse (environ \SI{-89}{\celsius}) : à \SI{25}{\celsius}, il est à l'état \textbf{gazeux}. (Le pentane est liquide, $\theta_{\text{éb}} \approx \SI{36}{\celsius}$ ; le décane est liquide, $\theta_{\text{éb}} \approx \SI{174}{\celsius}$.)
\end{enumerate}
\end{exoresolu}

\begin{attention}[Ne pas confondre évaporation et ébullition]
Un liquide s'\textbf{évapore} à toute température à sa surface (l'essence s'évapore à l'air libre), mais il ne \textbf{bout} qu'à une température précise, la température d'ébullition, où des bulles se forment dans toute la masse du liquide. Autre piège : ne dis pas « le gaz butane bout à \SI{-0,5}{\celsius}, donc il est liquide dans la bonbonne parce qu'il fait froid ». Dans la bonbonne, le butane est liquide parce qu'il est \textbf{comprimé} sous pression ; dès qu'on ouvre le robinet, la pression chute et il redevient gaz.
\end{attention}

\courssub{Conséquence : on peut séparer les constituants du pétrole}

Nous venons d'établir que chaque hydrocarbure du pétrole possède sa propre température d'ébullition, et que celle-ci dépend de la taille de la molécule. Voilà une information précieuse : si l'on chauffe du pétrole brut progressivement, les hydrocarbures légers (petites molécules) s'évaporeront en premier, puis les hydrocarbures de plus en plus lourds.

\begin{aretenir}[Principe de la séparation]
Les constituants d'un mélange liquide peuvent être \textbf{séparés par distillation} s'ils ont des \textbf{températures d'ébullition différentes}. Comme les hydrocarbures du pétrole bouillent à des températures d'autant plus élevées que leurs molécules sont grosses, on peut séparer le pétrole en plusieurs \textbf{fractions} par \textbf{distillation fractionnée}. C'est le principe du raffinage, que nous étudierons dans la section suivante.
\end{aretenir}

\begin{savaistu}[Pourquoi la SONARA chauffe le pétrole à plus de \SI{350}{\celsius} ?]
À la raffinerie de la SONARA (Société Nationale de Raffinage) à Limbé, le pétrole brut est chauffé à environ \SI{350}{\celsius} avant d'entrer dans la tour de distillation. À cette température, presque tous les hydrocarbures sont à l'état de vapeur, même les plus lourds comme l'hexadécane ($\theta_{\text{éb}} = \SI{287}{\celsius}$). Les vapeurs montent ensuite dans la tour, se refroidissent et se condensent à différents étages : chaque étage récupère une fraction différente. Sans la propriété que nous venons d'établir, aucune raffinerie ne pourrait fonctionner !
\end{savaistu}

\begin{exercice}[À toi de jouer]
1) Donne la formule brute de l'alcane possédant 7 atomes de carbone. Quel est son nom ?
2) Le dodécane \ce{C12H26} a une température d'ébullition d'environ \SI{216}{\celsius}. À l'aide de la courbe du cours, vérifie que cette valeur est cohérente avec la propriété établie.
3) Explique, en deux phrases, pourquoi on peut séparer l'essence du gasoil par distillation.
\end{exercice}

\begin{corrige}[Exercice n°1]
1) Pour $n = 7$ : $2n + 2 = 16$, donc la formule est \ce{C7H16} ; c'est l'\textbf{heptane}.
2) Le dodécane possède 12 atomes de carbone, donc sa température d'ébullition doit être comprise entre celle de l'octane ($n = 8$, \SI{126}{\celsius}) et celle de l'hexadécane ($n = 16$, \SI{287}{\celsius}). Or $126 < 216 < 287$ : la valeur est bien cohérente avec la courbe croissante.
3) L'essence est formée d'hydrocarbures légers (5 à 10 atomes de carbone) qui bouillent entre environ 40 et \SI{180}{\celsius}, alors que le gasoil contient des hydrocarbures plus lourds (12 à 20 atomes de carbone) bouillant entre environ 180 et \SI{350}{\celsius}. Leurs températures d'ébullition étant différentes, on les sépare en chauffant le pétrole : les vapeurs d'essence se condensent plus haut dans la tour que celles du gasoil.
\end{corrige}

\coursec{Le raffinage du pétrole : la distillation fractionnée}

Dans la section précédente, nous avons vu que le pétrole brut est un \cle{mélange} complexe de très nombreux \cle{hydrocarbures} : c'est un liquide noir, épais et visqueux, dont les molécules contiennent des chaînes carbonées de longueurs variées. Or, tel quel, ce brut sorti du puits est presque inutilisable : on ne peut pas mettre du pétrole brut dans le moteur d'une moto-taxi ni dans une lampe à kérosène. Il faut d'abord \emph{séparer} ses constituants. C'est toute la mission de la raffinerie, comme celle de la SONARA à Limbé. Comment séparer des liquides mélangés entre eux ? La réponse tient en une idée simple : ils ne bouillent pas tous à la même température.

\courssub{Qu'est-ce que le raffinage ?}

\begin{definition}[Raffinage]
Le \cle{raffinage} est l'ensemble des opérations qui transforment le \cle{pétrole brut} en produits directement utilisables : gaz, carburants (essence, kérosène, gasoil), huiles, fioul et bitume. Il se déroule dans une usine appelée \cle{raffinerie}.
\end{definition}

L'opération principale du raffinage est la \emph{distillation fractionnée}. Pour bien la comprendre, commençons par une expérience de classe, plus simple, qui en illustre le principe.

\begin{experience}[Distillation simple d'un mélange eau + encre]
\textbf{Matériel :} ballon, chauffe-ballon, thermomètre, réfrigérant à eau, erlenmeyer, mélange eau + quelques gouttes d'encre bleue, pierre ponce.\\
\textbf{Protocole :}
\begin{enumerate}
\item Verser le mélange coloré dans le ballon, avec quelques grains de pierre ponce (pour régulariser l'ébullition).
\item Chauffer doucement et observer le thermomètre.
\item Recueillir le liquide qui s'écoule dans l'erlenmeyer.
\end{enumerate}
\textbf{Observations :} la température se stabilise vers \SI{100}{\celsius} ; une vapeur incolore s'élève, se refroidit dans le réfrigérant et redevient liquide ; le liquide recueilli (le \cle{distillat}) est \emph{incolore} : c'est de l'eau pure. L'encre, qui ne se vaporise pas, reste dans le ballon.\\
\textbf{Interprétation :} l'eau bout et se vaporise, le colorant non. En condensant la vapeur, on récupère l'eau séparée du colorant. La distillation exploite la différence de \cle{température d'ébullition} des constituants d'un mélange.
\end{experience}

\begin{attention}[Sécurité pendant la distillation]
Ne \textbf{jamais} chauffer un récipient \emph{fermé} contenant un produit pétrolier (ou tout liquide inflammable) : la vapeur produite augmente la pression jusqu'à l'\cle{explosion}. Le montage doit toujours être ouvert côté sortie. De même, à la maison, ne jamais chauffer un bidon d'essence ou de gasoil bouché, et ne jamais approcher une flamme des vapeurs d'essence.
\end{attention}

\begin{popfigure}
\begin{center}
\begin{tikzpicture}[scale=0.9, every node/.style={font=\small}]
% Chauffe-ballon
\draw[fill=popmuted!40] (-0.9,0) rectangle (0.9,0.5);
\node at (0,0.25) {\footnotesize chauffe-ballon};
% Ballon
\draw[fill=popblueL] (0,1.6) circle (0.95);
\draw[fill=popblue!50] (-0.95,1.45) arc (200:340:0.98) -- cycle;
\draw (0,2.55) -- (0,3.2);
% Thermomètre
\draw[fill=white] (-0.12,2.4) rectangle (0.12,3.9);
\draw[fill=poporange] (-0.06,2.4) rectangle (0.06,3.3);
\node[right, font=\footnotesize] at (0.15,3.6) {thermomètre};
% Colonne de sortie
\draw (0.55,2.35) -- (2.2,3.0);
\draw (0.62,2.15) -- (2.27,2.8);
% Réfrigérant
\draw[fill=popblue!15] (2.2,2.55) -- (5.2,1.35) -- (5.45,1.75) -- (2.45,2.95) -- cycle;
\draw (2.45,2.75) -- (5.45,1.55);
\draw (2.2,2.75) -- (5.2,1.55);
\draw[->, popblue] (2.6,3.15) -- (2.6,2.95);
\draw[->, popblue] (4.9,1.15) -- (4.9,1.35);
\node[font=\footnotesize, popblue] at (2.3,3.35) {arrivée d'eau froide};
\node[font=\footnotesize, popblue] at (5.35,1.0) {sortie d'eau};
\node[font=\footnotesize] at (3.8,2.6) {réfrigérant};
% Erlenmeyer
\draw[fill=popblue!5] (5.0,0.1) -- (5.25,0.9) -- (5.25,1.25) -- (5.6,1.25) -- (5.6,0.9) -- (5.85,0.1) -- cycle;
\draw[fill=popblue!40] (5.12,0.12) -- (5.73,0.12) -- (5.6,0.5) -- (5.25,0.5) -- cycle;
\node[font=\footnotesize] at (5.42,-0.25) {distillat (eau pure)};
\node[font=\footnotesize] at (0,1.35) {mélange};
\end{tikzpicture}
\end{center}
\legende{Montage de distillation simple : le chauffe-ballon vaporise l'eau, le thermomètre contrôle la température, le réfrigérant refroidit la vapeur qui se condense, et le distillat est recueilli dans l'erlenmeyer. L'eau de refroidissement entre par le bas du réfrigérant et sort par le haut.}
\end{popfigure}

\courssub{Principe de la distillation fractionnée}

L'expérience précédente sépare \emph{deux} constituants. Le pétrole brut en contient des centaines ! On utilise alors une distillation \emph{fractionnée}, c'est-à-dire réalisée en plusieurs « fractions » successives.

\begin{definition}[Distillation fractionnée]
La \cle{distillation fractionnée} est une technique de séparation des constituants d'un mélange de liquides, fondée sur la différence de leurs \cle{températures d'ébullition}. Chaque constituant (ou groupe de constituants proches) se vaporise puis se condense à une température qui lui est propre : on obtient ainsi des \cle{fractions} distinctes.
\end{definition}

\begin{aretenir}[Idée-clé]
Plus une molécule d'hydrocarbure est \emph{courte} (peu d'atomes de carbone), plus elle est \emph{légère} et plus sa température d'ébullition est \emph{faible}. Plus la molécule est \emph{longue}, plus la fraction est \emph{lourde} et plus elle bout à \emph{haute} température.
\end{aretenir}

\courssub{Les étapes du raffinage dans la colonne à distillation}

\begin{methode}[Déroulement de la distillation fractionnée du pétrole]
\begin{enumerate}
\item \textbf{Chauffage :} le pétrole brut est chauffé dans un four jusqu'à environ \SI{400}{\celsius}. La plupart de ses constituants se vaporisent.
\item \textbf{Introduction dans la colonne :} le mélange de vapeurs et de liquide est injecté au bas d'une grande tour appelée \cle{colonne à distillation} (ou colonne à plateaux).
\item \textbf{Montée des vapeurs :} la colonne est de plus en plus \emph{froide} du bas vers le haut (environ \SI{400}{\celsius} en bas, moins de \SI{50}{\celsius} en haut). Les vapeurs montent en traversant des \cle{plateaux} percés.
\item \textbf{Condensation progressive :} chaque constituant se condense (redevient liquide) au niveau du plateau où la température est voisine de sa température d'ébullition. Les fractions \emph{lourdes} (forte température d'ébullition) se condensent \emph{en bas} ; les fractions \emph{légères} montent \emph{en haut}.
\item \textbf{Soutirage :} chaque fraction est recueillie (« soutirée ») à la hauteur correspondante. Le résidu solide, le \cle{bitume}, est récupéré au fond.
\end{enumerate}
\end{methode}

\begin{popfigure}
\begin{center}
\begin{tikzpicture}[scale=0.95, every node/.style={font=\small}]
% Four
\draw[fill=poporange!30] (-5.2,0.4) rectangle (-3.4,1.6);
\node[font=\footnotesize] at (-4.3,1.0) {four};
\node[font=\footnotesize] at (-4.3,0.15) {\SI{400}{\celsius}};
\draw[->, thick, popdark] (-6.4,1.0) -- (-5.2,1.0);
\node[left, font=\footnotesize] at (-6.4,1.0) {pétrole brut};
% Colonne
\draw[fill=popcream] (-2.0,-0.4) rectangle (0.4,7.2);
% Plateaux
\foreach \y in {0.3,1.2,2.1,3.0,3.9,4.8,5.7,6.6} {\draw[popmuted] (-2.0,\y) -- (0.4,\y);}
% Entrée du brut
\draw[->, thick, popdark] (-3.4,1.0) -- (-2.0,1.0);
% Flèche vapeurs montantes
\draw[->, very thick, poporange] (-0.8,0.6) -- (-0.8,6.9);
\node[rotate=90, font=\footnotesize, poporange] at (-1.15,3.7) {vapeurs montantes};
% Gradient de température
\draw[->, very thick, popblue] (0.0,6.9) -- (0.0,0.6);
\node[rotate=90, font=\footnotesize, popblue] at (0.35,3.7) {température croissante};
% Soutirages
\draw[->, thick, popgreen] (0.4,6.9) -- (1.8,6.9); \node[right, font=\footnotesize] at (1.8,6.9) {gaz (GPL) $< \SI{40}{\celsius}$};
\draw[->, thick, popgreen] (0.4,5.7) -- (1.8,5.7); \node[right, font=\footnotesize] at (1.8,5.7) {essences ($40$ à $\SI{180}{\celsius}$)};
\draw[->, thick, popgreen] (0.4,4.5) -- (1.8,4.5); \node[right, font=\footnotesize] at (1.8,4.5) {kérosène ($180$ à $\SI{250}{\celsius}$)};
\draw[->, thick, popgreen] (0.4,3.3) -- (1.8,3.3); \node[right, font=\footnotesize] at (1.8,3.3) {gasoil ($250$ à $\SI{350}{\celsius}$)};
\draw[->, thick, popgreen] (0.4,2.1) -- (1.8,2.1); \node[right, font=\footnotesize] at (1.8,2.1) {huiles lourdes, fioul};
\draw[->, thick, popdark] (0.4,-0.2) -- (1.8,-0.2); \node[right, font=\footnotesize] at (1.8,-0.2) {bitume (résidu)};
% Labels haut/bas
\node[font=\footnotesize\itshape, popblue] at (-0.8,7.5) {haut : froid, léger};
\node[font=\footnotesize\itshape, poporange] at (-0.8,-0.75) {bas : chaud, lourd};
\end{tikzpicture}
\end{center}
\legende{Colonne de distillation fractionnée : le brut chauffé à \SI{400}{\celsius} entre à la base ; les vapeurs montent et se condensent sur les plateaux d'autant plus haut qu'elles sont légères. La température décroît du bas vers le haut.}
\end{popfigure}

\courssub{Les fractions obtenues, du haut vers le bas}

\begin{center}
\begin{tabularx}{\linewidth}{|l|c|X|}\hline
\rowcolor{popblueL} \textbf{Fraction (haut $\to$ bas)} & \textbf{Ébullition} & \textbf{Utilisation principale} \\ \hline
Gaz de pétrole (GPL) & $< \SI{40}{\celsius}$ & gaz domestique (bonbonnes de gaz) \\ \hline
Essences & $40$ à $\SI{180}{\celsius}$ & carburant des motos et voitures \\ \hline
Kérosène & $180$ à $\SI{250}{\celsius}$ & carburant d'avion, lampes \\ \hline
Gasoil & $250$ à $\SI{350}{\celsius}$ & camions, groupes électrogènes \\ \hline
Huiles lourdes, fioul & $350$ à $\SI{400}{\celsius}$ & lubrifiants, combustible industriel \\ \hline
Bitume & résidu ($> \SI{400}{\celsius}$) & revêtement des routes \\ \hline
\end{tabularx}
\end{center}

\begin{methode}[Lire et interpréter le schéma d'une colonne de distillation]
\begin{enumerate}
\item Repérer le \textbf{sens du gradient de température} : la température \emph{décroît} du bas vers le haut (chaud en bas, froid en haut).
\item En déduire l'\textbf{ordre des fractions} : plus on monte, plus les molécules sont \emph{courtes} et la fraction \emph{légère}.
\item Vérifier la cohérence : les \emph{gaz} sortent tout en haut, le \emph{bitume} en bas.
\item Associer chaque soutirage à une utilisation concrète (GPL $\to$ bonbonne, essence $\to$ moto, bitume $\to$ route).
\end{enumerate}
\end{methode}

\begin{attention}[Erreur fréquente : inverser l'ordre des fractions]
Beaucoup d'élèves placent le bitume en haut et les gaz en bas. Retiens bien : les produits \textbf{légers} (gaz, essence), aux températures d'ébullition \emph{faibles}, montent et sortent \textbf{en haut} ; les produits \textbf{lourds} (fioul, bitume), aux températures d'ébullition \emph{élevées}, restent \textbf{en bas}. Astuce : « léger comme l'air, il monte ; lourd comme le goudron, il reste au fond ».
\end{attention}

\begin{savaistu}[Des tours géantes]
Une colonne de distillation industrielle peut mesurer plus de \SI{50}{m} de haut — l'équivalent d'un immeuble de 16 étages — et traiter des milliers de tonnes de brut par jour. À la SONARA de Limbé, ces colonnes ont longtemps fonctionné jour et nuit pour fournir le Cameroun en carburants.
\end{savaistu}

\begin{exemplebox}[Complément utile : le craquage (hors programme strict)]
La demande en essence est bien plus forte que la part d'essence naturellement présente dans le brut. Les raffineries pratiquent alors le \cle{craquage} : sous l'effet de la chaleur (et de catalyseurs), on « casse » les longues molécules des fractions \emph{lourdes} (fioul, huiles) pour obtenir des molécules \emph{courtes} d'essence. Le craquage transforme donc les fractions lourdes, moins demandées, en fractions légères très recherchées.
\end{exemplebox}

\begin{exoresolu}[Associer fraction et plateau]
Dans une colonne de distillation, on soutire quatre fractions aux températures de plateaux suivantes : \SI{60}{\celsius}, \SI{200}{\celsius}, \SI{300}{\celsius} et \SI{390}{\celsius}. Attribuer à chaque plateau la fraction recueillie parmi : bitume, gasoil, essence, kérosène. Préciser lequel est le plus haut.
\tcblower
\textbf{Solution.} La température décroît du bas vers le haut ; les fractions légères se condensent aux températures faibles, donc en haut.
\begin{itemize}
\item \SI{60}{\celsius} (plateau le plus haut) : \textbf{essence} (ébullition entre $40$ et $\SI{180}{\celsius}$) ;
\item \SI{200}{\celsius} : \textbf{kérosène} (entre $180$ et $\SI{250}{\celsius}$) ;
\item \SI{300}{\celsius} : \textbf{gasoil} (entre $250$ et $\SI{350}{\celsius}$) ;
\item \SI{390}{\celsius} (plateau le plus bas) : \textbf{bitume} (résidu lourd).
\end{itemize}
Le plateau le plus haut est celui à \SI{60}{\celsius}, où l'on recueille l'essence.
\end{exoresolu}

\begin{exercice}[À toi de jouer]
\begin{enumerate}
\item Pourquoi introduit-on le pétrole brut chauffé à la \emph{base} de la colonne et non au sommet ?
\item Classer ces fractions de la plus légère à la plus lourde : gasoil, GPL, bitume, essence, kérosène.
\item Expliquer, en deux phrases, pourquoi le craquage est économiquement intéressant pour une raffinerie.
\end{enumerate}
\end{exercice}

\begin{corrige}[Exercice « À toi de jouer »]
\begin{enumerate}
\item Le brut chauffé entre à la base, là où la température est la plus élevée, afin que toutes les fractions vaporisables se vaporisent ; les vapeurs montent ensuite et se condensent chacune à sa hauteur.
\item Du plus léger au plus lourd : GPL $<$ essence $<$ kérosène $<$ gasoil $<$ bitume.
\item Le craquage transforme les fractions lourdes, abondantes et peu demandées, en essence, carburant très recherché : la raffinerie augmente ainsi sa production d'essence et ses revenus.
\end{enumerate}
\end{corrige}

\begin{aretenir}[L'essentiel de la section]
\begin{itemize}
\item Le \cle{raffinage} transforme le pétrole brut en produits utilisables ; son opération clé est la \cle{distillation fractionnée}.
\item Principe : séparation selon les \cle{températures d'ébullition} ; le brut est chauffé vers \SI{400}{\celsius} puis introduit dans une colonne à plateaux.
\item La température \emph{décroît} du bas vers le haut ; les fractions \emph{légères} sortent en haut (GPL, essence), les \emph{lourdes} en bas (fioul, bitume).
\item Le \cle{craquage} (complément) transforme les fractions lourdes en fractions légères comme l'essence.
\end{itemize}
\end{aretenir}

\coursec{Les produits pétroliers et leurs utilisations}

Dans la section précédente, nous avons vu que le raffinage du pétrole repose sur la \cle{distillation fractionnée} : chauffé dans une tour de distillation, le pétrole brut se sépare en plusieurs \cle{fractions}, de la plus légère (les gaz) à la plus lourde (le bitume). Mais à quoi servent concrètement toutes ces fractions ? Pourquoi le chauffeur de ta moto-taxi met-il de l'essence et non du gasoil ? Pourquoi la bonbonne de gaz de la cuisine contient-elle du butane et non du fioul ? C'est ce que nous allons découvrir : chaque fraction pétrolière a une utilisation précise, qui dépend directement de sa composition, c'est-à-dire du \cle{nombre d'atomes de carbone} de ses molécules.

\courssub{Le principe général : la composition détermine l'usage}

\begin{aretenir}[L'idée directrice]
Plus les molécules d'une fraction pétrolière contiennent d'atomes de carbone, plus la fraction est \textbf{lourde} : elle s'évapore difficilement, s'enflamme moins facilement et bout à une température plus élevée.
\begin{itemize}
\item Peu d'atomes de carbone (1 à 4) : \textbf{gaz} à température ordinaire.
\item De 5 à 12 atomes de carbone environ : \textbf{liquides volatils} et très inflammables (essences).
\item De 12 à 25 atomes de carbone environ : \textbf{liquides épais} (kérosène, gasoil, huiles).
\item Plus de 25 atomes de carbone : \textbf{produits solides ou pâteux} (paraffines, bitume).
\end{itemize}
\end{aretenir}

\begin{methode}[Associer une fraction pétrolière à son utilisation]
Pour répondre à une question du type « quelle fraction utilise-t-on pour... ? », raisonne en trois étapes :
\begin{enumerate}
\item \textbf{Identifier le besoin} : faut-il un gaz (cuisine), un liquide très inflammable et volatile (moteur à essence), un liquide moins volatile mais énergétique (Diesel, avion), un liquide visqueux (lubrification) ou un solide (routes) ?
\item \textbf{Relier le besoin au nombre d'atomes de carbone} : gaz = 3 à 4 C ; essence = 5 à 10 C ; kérosène = 10 à 14 C ; gasoil = 14 à 20 C ; huiles et fiouls = 20 à 40 C ; bitume = plus de 40 C.
\item \textbf{Conclure en nommant la fraction} et en justifiant par sa propriété (volatilité, viscosité, inflammabilité).
\end{enumerate}
\end{methode}

\courssub{Les gaz de pétrole liquéfié (GPL) : butane et propane}

\begin{definition}[GPL]
Les \cle{gaz de pétrole liquéfié} (GPL) sont les fractions les plus légères du pétrole, formées d'hydrocarbures de 3 à 4 atomes de carbone : le \cle{propane} \ce{C3H8} et le \cle{butane} \ce{C4H10}. Ils sont gazeux à température ordinaire, mais on les \textbf{liquéfie par compression} pour les stocker dans des bonbonnes métalliques.
\end{definition}

Au Cameroun, le gaz butane est le \textbf{gaz domestique} par excellence : c'est lui qui alimente les réchauds et les cuisinières dans les cuisines de Douala, Yaoundé ou Bafoussam. Quand on ouvre le robinet de la bonbonne, la pression diminue : le butane liquide se transforme en gaz, qui brûle avec une flamme bleue selon la combustion complète :
\[ \ce{2C4H10 + 13O2 -> 8CO2 + 10H2O} \]
Vérifions l'équilibrage de cette équation : à gauche, $2 \times 4 = 8$ atomes de carbone, $2 \times 10 = 20$ atomes d'hydrogène et $13 \times 2 = 26$ atomes d'oxygène ; à droite, 8 atomes de carbone, $10 \times 2 = 20$ atomes d'hydrogène et $8 \times 2 + 10 = 26$ atomes d'oxygène. L'équation est bien équilibrée.

\begin{exemplebox}[Le gaz butane domestique]
Le gaz butane vendu en bonbonne est principalement du \ce{C4H10} : ses molécules contiennent \textbf{4 atomes de carbone}. C'est un gaz à température ordinaire, facile à liquéfier sous une pression modérée, ce qui permet de stocker une grande quantité d'énergie dans un petit volume. L'essence, elle, contient des hydrocarbures de \textbf{5 à 10 atomes de carbone} : c'est un liquide volatil, impossible à stocker dans une bonbonne comme le gaz.
\end{exemplebox}

\begin{attention}[Sécurité avec le gaz domestique]
Le butane est inodore : on lui ajoute une substance odorante pour détecter les fuites. En cas d'odeur de gaz : ne jamais actionner d'interrupteur ni de flamme, ouvrir les fenêtres, fermer le robinet de la bonbonne. Une fuite de gaz dans une cuisine fermée peut provoquer une explosion.
\end{attention}

\courssub{L'essence : carburant des moteurs à allumage commandé}

\begin{definition}[Essence]
L'\cle{essence} est un mélange d'hydrocarbures liquides contenant de 5 à 10 atomes de carbone environ. Très volatile et très inflammable, elle est le carburant des \cle{moteurs à allumage commandé} (moteurs à bougies) : motos, voitures, pirogues motorisées.
\end{definition}

Dans le moteur d'une moto-taxi, l'essence vaporisée est mélangée à l'air dans le cylindre, puis une \textbf{étincelle} produite par la bougie déclenche la combustion : c'est pour cela qu'on parle d'allumage \emph{commandé}. La volatilité de l'essence est justement ce qui permet de former facilement ce mélange air-carburant.

\courssub{Le kérosène : carburant d'aviation et éclairage}

\begin{definition}[Kérosène]
Le \cle{kérosène} est une fraction liquide contenant environ 10 à 14 atomes de carbone. Moins volatile que l'essence, il brûle de façon régulière et dégage beaucoup d'énergie.
\end{definition}

Le kérosène a deux grands usages :
\begin{itemize}
\item \textbf{Carburant d'aviation} : les réacteurs des avions (par exemple ceux qui décollent de l'aéroport de Douala) consomment du kérosène, car il reste liquide à basse température en altitude et fournit une combustion puissante et stable.
\item \textbf{Lampes et foyers} : dans certaines zones rurales non électrifiées, les lampes à pétrole et les foyers de cuisson utilisent le kérosène (souvent appelé « pétrole lampant »).
\end{itemize}

\courssub{Le gasoil : carburant des moteurs Diesel}

\begin{definition}[Gasoil]
Le \cle{gasoil} (ou gazole) est une fraction liquide contenant environ 14 à 20 atomes de carbone. Il est le carburant des \cle{moteurs Diesel} : camions, bus, tracteurs, trains et \textbf{groupes électrogènes}.
\end{definition}

Contrairement au moteur à essence, le moteur Diesel n'a \textbf{pas de bougie} : l'air est fortement comprimé dans le cylindre, ce qui l'échauffe, puis le gasoil injecté s'enflamme \textbf{spontanément} au contact de cet air chaud. On parle d'\emph{auto-inflammation}. C'est le gasoil qui fait tourner les groupes électrogènes des quartiers lors des délestages d'ENEO, et qui alimente les camions de transport sur l'axe Douala--Yaoundé.

\begin{attention}[Erreur fréquente : gasoil ou essence ?]
Ne confonds jamais les deux carburants :
\begin{itemize}
\item \textbf{Essence} $\rightarrow$ moteur \textbf{à allumage commandé} (bougie, étincelle) : motos, voitures.
\item \textbf{Gasoil} $\rightarrow$ moteur \textbf{Diesel} (auto-inflammation par compression, sans bougie) : camions, bus, groupes électrogènes.
\end{itemize}
Mettre de l'essence dans un moteur Diesel (ou l'inverse) détruit le moteur. Au BEPC, associer le bon carburant au bon moteur est un point de cours souvent demandé.
\end{attention}

\courssub{Les huiles, graisses et fiouls lourds}

\begin{definition}[Huiles et graisses lubrifiantes]
Les \cle{huiles} et \cle{graisses} sont des fractions lourdes (20 à 40 atomes de carbone environ), épaisses et visqueuses. Elles servent à la \cle{lubrification} des moteurs et des machines : elles forment un film entre les pièces métalliques en mouvement, ce qui réduit les frottements, l'usure et l'échauffement.
\end{definition}

Quand le mécanicien fait la « vidange » d'une moto ou d'une voiture, il remplace l'huile moteur usée. Sans huile, les pièces du moteur frotteraient directement l'une contre l'autre et le moteur « gripperait » en quelques minutes.

\begin{definition}[Fioul lourd]
Le \cle{fioul lourd} est une fraction très lourde et visqueuse. Il sert de \textbf{combustible} dans les centrales thermiques (production d'électricité) et dans les chaudières des grands navires.
\end{definition}

\courssub{Le bitume : le revêtement des routes}

\begin{definition}[Bitume]
Le \cle{bitume} est le \textbf{résidu le plus lourd} de la distillation du pétrole (plus de 40 atomes de carbone). Noir, solide à froid et pâteux à chaud, il sert de liant pour le \textbf{revêtement des routes} (goudronnage) et pour l'étanchéité des toitures.
\end{definition}

\begin{savaistu}[Le bitume des routes camerounaises]
Le bitume qui goudronne les routes camerounaises, comme l'axe Douala--Yaoundé, est le résidu le plus lourd de la distillation du pétrole : c'est ce qui reste au fond de la tour de distillation après le départ de toutes les autres fractions ! Chauffé pour devenir fluide, il est mélangé à du sable et à des gravillons pour former l'\textbf{enrobé} que l'on étale sur la chaussée. En refroidissant, il durcit et lie les gravillons entre eux. Rien ne se perd dans le raffinage : même le « fond de cuve » devient une richesse.
\end{savaistu}

\courssub{Complément : la pétrochimie}

\begin{definition}[Pétrochimie]
La \cle{pétrochimie} est l'industrie qui transforme certains produits pétroliers (notamment le naphta et les gaz) en nouvelles substances chimiques. Elle fabrique les \textbf{plastiques}, les \textbf{engrais}, les \textbf{médicaments}, les \textbf{textiles synthétiques}, les peintures et les détergents.
\end{definition}

Les sachets plastiques, les seaux, les tuyaux PVC, les tissus en polyester ou les engrais agricoles proviennent tous, à l'origine, du pétrole. Le pétrole n'est donc pas seulement un carburant : c'est aussi la matière première d'une grande partie des objets de notre vie quotidienne. C'est pourquoi la gestion des déchets plastiques est un enjeu majeur au Cameroun.

\courssub{Tableau récapitulatif des fractions pétrolières}

\begin{popfigure}
\begin{tikzpicture}[font=\small]
\fill[popblue] (0,0) rectangle (12,0.7);
\node[white,font=\bfseries] at (2.2,0.35) {Fraction};
\node[white,font=\bfseries] at (6.0,0.35) {Nombre de C (approx.)};
\node[white,font=\bfseries] at (9.8,0.35) {Utilisation courante};
\foreach \y/\fr/\nc/\use/\col in {
  -0.7/{GPL (butane, propane)}/3 \`a 4/Cuisine (gaz domestique)/popblueL,
  -1.4/Essences/5 \`a 10/{Motos, voitures (allumage command\'e)}/white,
  -2.1/K\'eros\`ene/10 \`a 14/{Avions, lampes, foyers}/popblueL,
  -2.8/Gasoil/14 \`a 20/{Camions, bus, groupes \'electrog\`enes}/white,
  -3.5/Huiles et graisses/20 \`a 40/Lubrification des moteurs/popblueL,
  -4.2/Fioul lourd/25 \`a 40/{Centrales thermiques, navires}/white,
  -4.9/Bitume/plus de 40/Rev\^etement des routes/popblueL}{
  \fill[\col] (0,\y) rectangle (12,\y+0.7);
  \node[anchor=west] at (0.2,\y+0.35) {\fr};
  \node at (6.0,\y+0.35) {\nc};
  \node[anchor=west] at (7.6,\y+0.35) {\use};
}
\draw[popdark,thick] (0,0.7) rectangle (12,-4.9);
\draw[popdark] (4.4,0.7) -- (4.4,-4.9);
\draw[popdark] (7.5,0.7) -- (7.5,-4.9);
\foreach \y in {-0.7,-1.4,-2.1,-2.8,-3.5,-4.2}{\draw[popline] (0,\y) -- (12,\y);}
\end{tikzpicture}
\legende{Les principales fractions du pétrole, leur composition approximative en nombre d'atomes de carbone et leur utilisation courante. De haut en bas, les fractions sont de plus en plus lourdes.}
\end{popfigure}

\begin{popfigure}
\begin{center}
\begin{tikzpicture}
\begin{axis}[
  ybar, bar width=0.6cm,
  width=0.9\linewidth, height=6cm,
  ymin=0, ymax=70,
  ylabel={Part approximative (\%)},
  symbolic x coords={Carburants,Chauffage,Petrochimie,Autres},
  xtick=data,
  xticklabels={Carburants,Chauffage,P\'etrochimie,Autres},
  x tick label style={font=\small},
  nodes near coords,
  every node near coord/.append style={font=\small},
  axis lines=left,
  ymajorgrids=true, grid style={popline!50},
]
\addplot[fill=poporange, draw=popdark] coordinates {(Carburants,55) (Chauffage,15) (Petrochimie,20) (Autres,10)};
\end{axis}
\end{tikzpicture}
\end{center}
\legende{Répartition approximative des débouchés du pétrole dans le monde : la majeure partie devient des carburants (transports), le reste sert au chauffage, à la pétrochimie (plastiques, engrais, médicaments) et à d'autres usages (lubrifiants, bitume).}
\end{popfigure}

\begin{exoresolu}[Choisir la bonne fraction]
Énoncé : Un transporteur camerounais possède un camion à moteur Diesel, une moto à moteur à allumage commandé et une cuisinière à gaz. Pour chaque engin, indique la fraction pétrolière utilisée et justifie ta réponse par la composition de cette fraction.
\tcblower
Solution détaillée :
\begin{enumerate}
\item \textbf{Camion à moteur Diesel} : il utilise le \textbf{gasoil}, fraction contenant environ 14 à 20 atomes de carbone. Le gasoil s'enflamme spontanément dans l'air fortement comprimé du cylindre (auto-inflammation), ce qui convient au moteur Diesel sans bougie.
\item \textbf{Moto à allumage commandé} : elle utilise l'\textbf{essence}, mélange d'hydrocarbures de 5 à 10 atomes de carbone. Très volatile, l'essence se vaporise et se mélange à l'air ; une étincelle de la bougie déclenche la combustion.
\item \textbf{Cuisinière à gaz} : elle utilise le \textbf{butane} (\ce{C4H10}), un GPL de 4 atomes de carbone, gazeux à température ordinaire et stocké liquide sous pression dans la bonbonne.
\end{enumerate}
Conclusion : à chaque usage correspond une fraction, choisie en fonction du nombre d'atomes de carbone de ses molécules, donc de sa volatilité et de son mode de combustion.
\end{exoresolu}

\begin{exercice}[À toi de jouer]
\begin{enumerate}
\item Pourquoi stocke-t-on le butane sous pression dans une bonbonne alors que l'essence se conserve dans un simple bidon ?
\item Un élève affirme : « Le gasoil et l'essence, c'est pareil, ce sont des carburants. » Corrige cette erreur en deux phrases.
\item Classe ces fractions de la plus légère à la plus lourde : bitume, essence, gasoil, GPL, kérosène.
\item Cite deux objets de ta vie quotidienne fabriqués grâce à la pétrochimie.
\end{enumerate}
\end{exercice}

\begin{corrige}[Exercice : À toi de jouer]
\begin{enumerate}
\item Le butane est un gaz à température ordinaire : on le comprime pour le liquéfier et stocker beaucoup de gaz dans un petit volume. L'essence est déjà liquide à température ordinaire : un simple bidon fermé suffit.
\item L'essence alimente les moteurs à allumage commandé (avec bougie), comme ceux des motos. Le gasoil alimente les moteurs Diesel (sans bougie, auto-inflammation par compression), comme ceux des camions ; ce sont deux carburants différents, non interchangeables.
\item GPL $<$ essence $<$ kérosène $<$ gasoil $<$ bitume.
\item Exemples : sachets plastiques, seaux en plastique, tissus en polyester, engrais, certains médicaments.
\end{enumerate}
\end{corrige}

\begin{aretenir}[L'essentiel de la section]
\begin{itemize}
\item Chaque fraction pétrolière a une utilisation qui dépend du \textbf{nombre d'atomes de carbone} de ses molécules : plus il est grand, plus la fraction est lourde et visqueuse.
\item \textbf{GPL} (butane, propane) : gaz domestique pour la cuisine. \textbf{Essence} : moteurs à allumage commandé (motos, voitures). \textbf{Kérosène} : avions, lampes. \textbf{Gasoil} : moteurs Diesel (camions, groupes électrogènes). \textbf{Huiles} : lubrification. \textbf{Fioul lourd} : centrales thermiques, navires. \textbf{Bitume} : revêtement des routes.
\item La \textbf{pétrochimie} transforme les produits pétroliers en plastiques, engrais, médicaments et textiles.
\end{itemize}
\end{aretenir}

\coursec{Pétrole et pollution : impacts et lutte}

Dans la section précédente, nous avons découvert les nombreux produits issus du raffinage du pétrole : essence, gasoil, kérosène, fioul, bitume, plastiques... Ces produits rendent d'immenses services : ils font rouler les moto-taxis, alimentent les groupes électrogènes quand ENEO déleste, chauffent les marmites. Mais chaque médaille a son revers. Brûler des carburants, répandre des huiles, jeter des sachets plastiques dans la rue : tout cela dégrade notre air, nos sols et nos eaux. Dans cette dernière section, nous allons identifier les grandes pollutions liées au pétrole, comprendre leurs mécanismes, puis apprendre les mesures de lutte et les gestes citoyens qui protègent notre environnement.

\courssub{La pollution de l'air par les produits pétroliers}

\textbf{Observation de la vie courante.} Derrière un vieux camion qui grimpe une côte à Yaoundé, on voit parfois une épaisse fumée noire. Près d'un groupe électrogène mal réglé, l'odeur est âcre et pique les yeux. Que contiennent ces fumées ? Sont-elles dangereuses ? C'est toute la question de cette sous-section.

\textbf{Hypothèse.} La fumée noire et les odeurs désagréables traduisent une mauvaise combustion du carburant, qui libère des gaz nocifs en plus des gaz normaux.

\begin{experience}[Fumées d'une bougie : combustion complète ou incomplète ?]
\textbf{Matériel :} une bougie, une soucoupe en verre ou en métal propre, une boîte d'allumettes.

\textbf{Protocole :}
\begin{enumerate}
\item Allume la bougie et tiens la soucoupe \emph{au-dessus} de la flamme, dans la zone claire, quelques secondes. Observe le dépôt éventuel.
\item Recommence en plaçant la soucoupe \emph{dans} la flamme, de façon à gêner l'arrivée d'air. Observe.
\end{enumerate}

\textbf{Observations :} dans le premier cas, le dépôt est faible ; dans le second cas, la soucoupe se couvre d'un dépôt \cle{noir} et poudreux : c'est du \cle{carbone} (suie).

\textbf{Interprétation :} quand l'air (donc le dioxygène) manque, la combustion de la cire — un hydrocarbure, comme les carburants — est \emph{incomplète} : au lieu de produire seulement du dioxyde de carbone et de l'eau, elle produit aussi du carbone solide (fumées noires) et un gaz invisible et mortel, le monoxyde de carbone \ce{CO}.

\textbf{Conclusion :} la qualité de la combustion dépend de la quantité de dioxygène disponible. Une combustion incomplète est beaucoup plus polluante.
\end{experience}

\begin{definition}[Combustion complète et combustion incomplète (rappel)]
La \cle{combustion complète} d'un hydrocarbure se déroule avec une quantité de dioxygène suffisante : elle produit uniquement du \cle{dioxyde de carbone} \ce{CO2} et de l'\cle{eau} \ce{H2O}.

La \cle{combustion incomplète} se déroule avec une quantité de dioxygène insuffisante : elle produit, en plus du dioxyde de carbone et de l'eau, du \cle{carbone} \ce{C} (suies, fumées noires) et du \cle{monoxyde de carbone} \ce{CO}, gaz toxique.
\end{definition}

\begin{exemplebox}[Équations de combustion du méthane]
Le méthane \ce{CH4} est l'hydrocarbure le plus simple (c'est le constituant principal du gaz naturel).

Combustion \emph{complète} (dioxygène en excès) :
\[ \ce{CH4 + 2O2 -> CO2 + 2H2O} \]

Combustion \emph{incomplète} (dioxygène insuffisant) : il se forme du monoxyde de carbone ou du carbone :
\[ \ce{2CH4 + 3O2 -> 2CO + 4H2O} \qquad \ce{CH4 + O2 -> C + 2H2O} \]
Vérifie sur chaque équation que le nombre d'atomes de chaque élément est le même avant et après la flèche : c'est la conservation des atomes.
\end{exemplebox}

\begin{aretenir}[Les quatre grandes pollutions de l'air liées au pétrole]
\begin{itemize}
\item Le \cle{dioxyde de carbone} \ce{CO2}, produit par toute combustion complète, s'accumule dans l'atmosphère et renforce l'\cle{effet de serre}, ce qui provoque le réchauffement climatique.
\item Les \cle{oxydes de soufre} (venant du soufre présent dans le fioul et le gasoil) et les \cle{oxydes d'azote} (formés à haute température dans les moteurs) se transforment dans l'atmosphère en acides : ils provoquent les \cle{pluies acides}, qui acidifient les sols et les lacs, abîment les forêts et rongent les toitures en tôle.
\item Les \cle{fumées noires} (particules de carbone) et le \cle{monoxyde de carbone} \ce{CO} proviennent des combustions \emph{incomplètes} : moteurs mal réglés, feux de déchets, groupes électrogènes encrassés.
\end{itemize}
\end{aretenir}

\begin{attention}[Erreur fréquente : « seul le \ce{CO2} pollue »]
Beaucoup d'élèves croient que la seule pollution des carburants est le \ce{CO2}. C'est faux ! Le \ce{CO2} est le produit \emph{normal} d'une combustion complète ; son problème est son accumulation dans l'atmosphère (effet de serre). Le gaz le plus immédiatement dangereux est le \cle{monoxyde de carbone} \ce{CO} : invisible, sans odeur, il se fixe sur le sang à la place du dioxygène et peut tuer en quelques minutes dans une pièce mal ventilée. Il provient des \emph{combustions incomplètes}. C'est pourquoi on ne fait jamais tourner un groupe électrogène dans une pièce fermée.
\end{attention}

\courssub{La pollution des sols et des eaux}

\textbf{Situation concrète.} En mer, un pétrolier accidenté répand des milliers de tonnes de pétrole brut qui souillent les plages et étouffent oiseaux et poissons. Sur terre, un oléoduc qui fuit imbibe les sols de brut. Dans nos quartiers, certains mécaniciens versent l'huile de vidange des moteurs directement sur le sol ou dans les caniveaux : à la première pluie, cette huile rejoint les ruisseaux et les nappes phréatiques.

\begin{definition}[Marée noire]
Une \cle{marée noire} est une pollution marine causée par un déversement accidentel de pétrole brut (ou de produits pétroliers) en mer, à la suite d'un naufrage de pétrolier, d'une fuite de plateforme ou d'une rupture d'oléoduc. Le pétrole, moins dense que l'eau et non miscible à elle, s'étale en nappe à la surface, étouffe la faune et la flore marines et souille les côtes.
\end{definition}

\begin{aretenir}[Les trois sources principales de pollution des sols et des eaux]
\begin{itemize}
\item Les \cle{marées noires} en mer ;
\item les \cle{fuites d'oléoducs} et les déversements lors de l'extraction ou du transport, qui contaminent les terres agricoles ;
\item les \cle{vidanges sauvages} des huiles usagées (moteurs, groupes électrogènes) jetées dans la nature ou les caniveaux.
\end{itemize}
\end{aretenir}

\begin{savaistu}[Un litre d'huile, un million de litres d'eau pollués]
Un seul litre d'huile de vidange jeté dans la nature peut polluer jusqu'à \textbf{un million de litres d'eau} : il forme à la surface une pellicule qui empêche le dioxygène de l'air de se dissoudre dans l'eau, ce qui asphyxie les poissons. Or les huiles usagées peuvent être \emph{collectées, régénérées et réutilisées} ! Voilà pourquoi il faut rapporter son huile de vidange chez un garagiste sérieux ou à un point de collecte, et jamais la verser au sol.
\end{savaistu}

\courssub{La pollution par les plastiques issus de la pétrochimie}

Les plastiques (sachets, bouteilles, emballages) sont fabriqués à partir du pétrole par la \cle{pétrochimie}. Leur défaut majeur : ils ne sont \emph{pas biodégradables}, c'est-à-dire que les micro-organismes du sol ne parviennent pas à les décomposer. Un sachet plastique jeté dans la rue peut mettre plusieurs \emph{centaines d'années} à disparaître. Conséquences visibles partout au Cameroun :
\begin{itemize}
\item accumulation des déchets dans les rues, les champs et les cours d'eau ;
\item \cle{colmatage des caniveaux}, qui provoque des inondations en saison des pluies et favorise la prolifération des moustiques (paludisme) ;
\item brûlage sauvage des plastiques, qui dégage des fumées toxiques ;
\item ingestion de plastiques par les animaux (chèvres, poissons), qui les tuent ou contaminent notre alimentation.
\end{itemize}

\courssub{Les mesures de lutte contre les pollutions pétrolières}

\begin{methode}[Proposer une mesure de lutte adaptée à chaque pollution]
Face à un problème de pollution, raisonne en quatre étapes :
\begin{enumerate}
\item \textbf{Identifier} le milieu pollué : air, sol, eau ?
\item \textbf{Identifier} le polluant : \ce{CO2}, \ce{CO}, suies, pluies acides, huile, plastique ?
\item \textbf{Identifier} la cause : combustion incomplète, rejet, fuite, déchet jeté ?
\item \textbf{Associer} la mesure qui supprime ou réduit cette cause.
\end{enumerate}
Exemples d'associations : fumées noires et \ce{CO} $\rightarrow$ régler le moteur ; oxydes d'azote $\rightarrow$ pot catalytique ; huile usagée $\rightarrow$ recyclage ; sachets $\rightarrow$ tri et recyclage des plastiques.
\end{methode}

\begin{aretenir}[Les grandes mesures de lutte]
\begin{itemize}
\item \textbf{Moteurs bien réglés} (voitures, motos, groupes électrogènes) : la combustion devient complète, on élimine fumées noires et \ce{CO}, et on consomme moins de carburant.
\item \textbf{Pots catalytiques} sur les véhicules : ils transforment une partie des gaz toxiques (\ce{CO}, oxydes d'azote) en gaz moins nocifs (\ce{CO2}, \ce{N2}).
\item \textbf{Recyclage} des huiles usagées (régénération) et des plastiques (fonte et transformation en nouveaux objets : pavés, ardoises, tuyaux).
\item \textbf{Législation et contrôles} : traitement des eaux de ballast des pétroliers, interdiction des dégazages en mer, amendes contre les déversements sauvages.
\item \textbf{Économies d'énergie} : éteindre les appareils inutiles, entretenir les moteurs, privilégier les transports en commun.
\item \textbf{Énergies renouvelables} (solaire, hydraulique, éolien, biomasse) : solution à long terme pour remplacer progressivement les énergies fossiles. Le Cameroun, avec son fort potentiel hydroélectrique et solaire, y a tout intérêt.
\end{itemize}
\end{aretenir}

\begin{popfigure}
\begin{tikzpicture}[scale=1, every node/.style={align=center}]
\node[draw=popblue, fill=popblueL, rounded corners, minimum width=2.6cm, minimum height=1cm, font=\small\bfseries, align=center] (ext) at (0,0){Extraction\\(gisements)};
\node[draw=poporange, fill=poporangeL, rounded corners, minimum width=2.6cm, minimum height=1cm, font=\small\bfseries, align=center] (raf) at (4.2,0){Raffinage\\(distillation)};
\node[draw=popgreen, fill=popgreenL, rounded corners, minimum width=2.6cm, minimum height=1cm, font=\small\bfseries, align=center] (uti) at (8.4,0){Utilisations\\(carburants, plastiques)};
\node[draw=poppink, fill=poppinkL, rounded corners, minimum width=3cm, minimum height=1cm, font=\small\bfseries, align=center] (pol) at (8.4,-2.6){Pollutions\\(air, sols, eaux)};
\node[draw=poppurple, fill=poppurpleL, rounded corners, minimum width=3cm, minimum height=1cm, font=\small\bfseries, align=center] (lut) at (4.2,-2.6){Mesures de lutte\\(recyclage, réglage)};
\node[draw=popgold, fill=popgoldL, rounded corners, minimum width=3cm, minimum height=1cm, font=\small\bfseries, align=center] (eco) at (0,-2.6){Économies d'énergie\\énergies renouvelables};
\draw[-{Stealth}, very thick, popdark] (ext) -- (raf);
\draw[-{Stealth}, very thick, popdark] (raf) -- (uti);
\draw[-{Stealth}, very thick, popdark] (uti) -- (pol);
\draw[-{Stealth}, very thick, popdark] (pol) -- (lut);
\draw[-{Stealth}, very thick, popdark] (lut) -- (eco);
\draw[-{Stealth}, very thick, popgreen] (eco) -- (ext) node[midway, left, font=\scriptsize\itshape, text=popdark]{consommer moins};
\end{tikzpicture}
\legende{La boucle du pétrole : de l'extraction aux mesures de lutte. Les économies d'énergie et les énergies renouvelables permettent de réduire la demande en pétrole et donc toutes les pollutions associées.}
\end{popfigure}

\courssub{Sécurité : stockage et manipulation des produits pétroliers}

Les produits pétroliers (essence, gasoil, kérosène) sont inflammables et leurs vapeurs toxiques. Quelques règles simples sauvent des vies :
\begin{itemize}
\item stocker les carburants \emph{loin des flammes, des cuisines et des habitations}, dans un local \cle{ventilé} ;
\item utiliser des \cle{bidons fermés et étiquetés} (jamais de bouteilles de boisson : risque d'ingestion par un enfant !) ;
\item ne \emph{jamais} manipuler d'hydrocarbures près d'un feu, d'une cigarette ou d'un appareil électrique qui produit des étincelles ;
\item ne jamais siphonner de l'essence avec la bouche.
\end{itemize}

\begin{attention}[Sécurité : les vapeurs d'essence, un danger invisible]
Les vapeurs d'essence sont \emph{invisibles} et \emph{plus denses que l'air} : elles s'accumulent au ras du sol et peuvent s'enflammer à plusieurs mètres du bidon, puis ramener la flamme jusqu'au dépôt. Règle absolue : \textbf{aucune flamme, aucune étincelle, aucune cigarette près d'un dépôt de carburant}. Une simple étincelle d'interrupteur ou d'appareil électrique peut suffire à déclencher une explosion.
\end{attention}

\begin{popfigure}
\begin{tikzpicture}[scale=1]
% Inflammable : flamme dans losange
\begin{scope}[shift={(0,0)}]
\draw[fill=poppinkL, draw=poppink, very thick, rotate=45] (-0.9,-0.9) rectangle (0.9,0.9);
\draw[fill=poporange, draw=popdark] (0,-0.35) .. controls (0.35,0.05) and (0.3,0.45) .. (0,0.75) .. controls (-0.3,0.45) and (-0.35,0.05) .. (0,-0.35);
\draw[fill=popgold] (0,-0.15) .. controls (0.15,0.1) and (0.12,0.3) .. (0,0.45) .. controls (-0.12,0.3) and (-0.15,0.1) .. (0,-0.15);
\node[font=\scriptsize\bfseries, text=popdark] at (0,-1.55) {Inflammable};
\end{scope}
% Explosion : étoile éclatée
\begin{scope}[shift={(4,0)}]
\draw[fill=popgoldL, draw=popgold, very thick, rotate=45] (-0.9,-0.9) rectangle (0.9,0.9);
\foreach \a in {0,45,...,315} \draw[very thick, popdark] (0,0) -- (\a:0.55);
\draw[fill=poporange, draw=popdark] (0,0) circle (0.22);
\node[font=\scriptsize\bfseries, text=popdark] at (0,-1.55) {Danger d'explosion};
\end{scope}
% Toxique : tête de mort stylisée
\begin{scope}[shift={(8,0)}]
\draw[fill=poppurpleL, draw=poppurple, very thick, rotate=45] (-0.9,-0.9) rectangle (0.9,0.9);
\draw[fill=white, draw=popdark, thick] (0,0.12) circle (0.4);
\fill[popdark] (-0.14,0.2) circle (0.07);
\fill[popdark] (0.14,0.2) circle (0.07);
\draw[draw=popdark, thick] (-0.15,-0.12) -- (-0.15,-0.28);
\draw[draw=popdark, thick] (0,-0.12) -- (0,-0.28);
\draw[draw=popdark, thick] (0.15,-0.12) -- (0.15,-0.28);
\draw[draw=popdark, thick, rounded corners] (-0.28,-0.14) -- (-0.28,-0.3) -- (0.28,-0.3) -- (0.28,-0.14);
\node[font=\scriptsize\bfseries, text=popdark] at (0,-1.55) {Produit toxique};
\end{scope}
\end{tikzpicture}
\legende{Pictogrammes de sécurité à connaître : produit inflammable (flamme), danger d'explosion (éclat), produit toxique (tête de mort). Ils figurent sur les bidons et fûts de produits pétroliers.}
\end{popfigure}

\courssub{Responsabilité citoyenne : les gestes qui protègent}

La lutte contre la pollution n'est pas réservée aux gouvernements et aux industriels. Chacun de nous peut agir :
\begin{itemize}
\item ne \emph{jamais} jeter l'huile de vidange dans la nature, les caniveaux ou les champs : la rapporter à un point de collecte ;
\item limiter les sachets plastiques (préférer les sacs réutilisables), trier et recycler les bouteilles ;
\item ne \emph{jamais} brûler les plastiques à l'air libre ;
\item entretenir le moteur de la moto ou du groupe électrogène de la famille ;
\item éteindre les lampes et appareils inutiles pour économiser l'énergie.
\end{itemize}

\begin{exoresolu}[Associer pollution et mesure de lutte]
\textbf{Énoncé.} Dans un quartier de Douala, on observe : (a) des fumées noires derrière les vieux camions ; (b) de l'huile de vidange versée dans un caniveau ; (c) des sachets plastiques qui bouchent les caniveaux et provoquent des inondations. Pour chaque cas, identifie le type de pollution et propose une mesure de lutte adaptée.
\tcblower
\textbf{Solution détaillée.} On applique la méthode en quatre étapes (milieu, polluant, cause, mesure).

\textbf{Cas (a) :} milieu pollué : l'\emph{air} ; polluants : suies (fumées noires) et monoxyde de carbone \ce{CO} ; cause : \emph{combustion incomplète} du gasoil dans des moteurs mal réglés ; mesure : \textbf{régler et entretenir les moteurs} (et équiper les véhicules de pots catalytiques).

\textbf{Cas (b) :} milieux pollués : les \emph{sols et les eaux} ; polluant : huile usagée (un litre peut polluer jusqu'à un million de litres d'eau) ; cause : vidange sauvage ; mesure : \textbf{collecte et recyclage des huiles usagées}, avec contrôles et amendes.

\textbf{Cas (c) :} milieux pollués : les \emph{sols et les eaux} ; polluant : déchets plastiques non biodégradables ; cause : jets sauvages et absence de tri ; mesure : \textbf{tri, collecte et recyclage des plastiques}, et réduction de l'usage des sachets (sacs réutilisables).
\end{exoresolu}

\begin{exercice}[À toi de jouer]
\begin{enumerate}
\item Définis une marée noire et cite deux de ses conséquences sur le milieu marin.
\item Pourquoi ne doit-on jamais faire fonctionner un groupe électrogène dans une pièce fermée ? Nomme le gaz responsable du danger.
\item Cite trois mesures permettant de lutter contre la pollution de l'air par les véhicules.
\item Explique pourquoi les vapeurs d'essence sont particulièrement dangereuses près d'un dépôt de carburant.
\end{enumerate}
\end{exercice}

\begin{corrige}[Exercice : À toi de jouer]
\begin{enumerate}
\item Une marée noire est une pollution marine due à un déversement accidentel de pétrole en mer. Conséquences : le pétrole étouffe poissons, oiseaux et plancton, et souille les plages et les côtes.
\item Dans une pièce fermée, la combustion devient incomplète (manque de dioxygène) et produit du \cle{monoxyde de carbone} \ce{CO}, gaz invisible, sans odeur et mortel, qui se fixe sur le sang à la place du dioxygène.
\item Régler les moteurs (combustion complète), équiper les véhicules de pots catalytiques, développer les transports en commun et les énergies renouvelables.
\item Les vapeurs d'essence sont invisibles, plus denses que l'air et très inflammables : elles s'accumulent au sol et peuvent s'enflammer à distance, ramenant le feu jusqu'au dépôt. D'où l'interdiction de toute flamme ou étincelle à proximité.
\end{enumerate}
\end{corrige}

\begin{aretenir}[L'essentiel de la section]
\begin{itemize}
\item La combustion des produits pétroliers pollue l'air : \ce{CO2} (effet de serre), oxydes de soufre et d'azote (pluies acides), suies et \ce{CO} (combustion incomplète, gaz mortel).
\item Les sols et les eaux sont pollués par les marées noires, les fuites d'oléoducs et les vidanges sauvages d'huiles usagées.
\item Les plastiques, non biodégradables, s'accumulent dans l'environnement et bouchent les caniveaux.
\item Les solutions : moteurs réglés, pots catalytiques, recyclage des huiles et des plastiques, législation, économies d'énergie et énergies renouvelables.
\item Sécurité : stocker les carburants loin des flammes, dans des bidons fermés et étiquetés, dans des lieux ventilés ; aucune étincelle près d'un dépôt.
\end{itemize}
\end{aretenir}

\newpage

\coursec{Activité d'intégration}

\courssub{Objectifs de l'activité}

À la fin de cette activité, tu seras capable de :

\begin{itemize}
\item exploiter un tableau de données expérimentales (températures d'ébullition des alcanes) et tracer une courbe simple ;
\item établir le lien entre le nombre d'atomes de carbone d'un alcane et sa température d'ébullition ;
\item légender le schéma d'une colonne de distillation fractionnée en plaçant les fractions dans le bon ordre ;
\item identifier, dans la vie courante camerounaise, la fraction pétrolière utilisée par les motos-taxis et celle des cuisinières à gaz ;
\item proposer des mesures concrètes de lutte contre la pollution et des consignes de sécurité adaptées à un dépôt de carburant.
\end{itemize}

\courssub{Situation}

\textbf{De la nappe de brut de Limbé à la moto de quartier : le parcours du pétrole.}

La raffinerie de la SONARA, située à Limbé dans la région du Sud-Ouest, traite le pétrole brut extrait des gisements camerounais. Ce brut est un mélange complexe d'hydrocarbures qu'il faut séparer pour obtenir des produits utilisables : gaz domestique, essence pour les motos-taxis, gasoil pour les camions et les groupes électrogènes, fioul, bitume pour les routes.

Ton professeur de PCT te remet trois documents.

\textbf{Document 1 : températures d'ébullition de quelques alcanes (sous la pression atmosphérique normale).}

\begin{center}
\begin{tabularx}{\linewidth}{|l|X|X|X|}
\hline
\rowcolor{popblueL} \textbf{Alcane} & \textbf{Formule brute} & \textbf{Nombre d'atomes de C} & \textbf{Température d'ébullition (\SI{}{\celsius})} \\
\hline
Méthane & \ce{CH4} & 1 & $-162$ \\
\hline
Butane & \ce{C4H10} & 4 & $-0,5$ \\
\hline
Pentane & \ce{C5H12} & 5 & $36$ \\
\hline
Hexane & \ce{C6H14} & 6 & $69$ \\
\hline
Octane & \ce{C8H18} & 8 & $126$ \\
\hline
Décane & \ce{C10H22} & 10 & $174$ \\
\hline
Dodécane & \ce{C12H26} & 12 & $216$ \\
\hline
Hexadécane & \ce{C16H34} & 16 & $287$ \\
\hline
\end{tabularx}
\end{center}

\textbf{Document 2 : schéma muet d'une colonne de distillation fractionnée.} La colonne est chauffée à environ \SI{400}{\celsius} à sa base. La température diminue de bas en haut. Six soutirages (sorties latérales) sont numérotés de 1 (en haut, zone la plus froide) à 6 (en bas, zone la plus chaude), plus un résidu recueilli au fond. Les fractions à placer sont : gaz de pétrole (butane, propane), essence, kérosène, gasoil, fioul, bitume (résidu).

\textbf{Document 3 : extrait d'un article de presse.} « Quartier Ndogpassi, Douala. Une fuite d'huile provenant d'un dépôt clandestin de carburant a souillé les caniveaux et menace la nappe phréatique. Par ailleurs, de plus en plus de familles cuisinent au gaz butane en bonbonne, tandis que les motos-taxis consomment chaque jour des milliers de litres d'essence. Les autorités rappellent que le stockage et la vente de carburant doivent respecter des règles strictes de sécurité. »

\courssub{Consignes}

Travail en groupes de quatre élèves, puis restitution orale et correction collective.

\begin{enumerate}
\item À partir du document 1, trace sur papier millimétré la courbe donnant la température d'ébullition en fonction du nombre d'atomes de carbone. Axe des abscisses : nombre d'atomes de C (1 cm pour 1 atome) ; axe des ordonnées : température (1 cm pour \SI{20}{\celsius}).
\item Décris l'évolution de la température d'ébullition quand le nombre d'atomes de carbone augmente. Formule une conclusion générale.
\item Explique, en t'appuyant sur cette conclusion, pourquoi la distillation fractionnée permet de séparer les constituants du pétrole brut.
\item Légende le schéma muet du document 2 : place les six fractions (gaz de pétrole, essence, kérosène, gasoil, fioul, bitume) aux soutirages 1 à 6 et au fond de la colonne. Justifie ton choix.
\item À l'aide du document 3 et de tes connaissances, indique : (a) quelle fraction alimente les motos-taxis ; (b) quelle fraction contient le gaz butane des cuisinières ; (c) où ces deux fractions sont soutirées dans la colonne.
\item Rédige un court avis citoyen (10 à 15 lignes) adressé au chef de quartier de Ndogpassi, proposant trois mesures concrètes de prévention de la pollution par les hydrocarbures et deux consignes de sécurité pour un dépôt de carburant de quartier.
\end{enumerate}

\courssub{Ressources à mobiliser}

\begin{itemize}
\item Le pétrole brut est un \cle{mélange} d'hydrocarbures, principalement des alcanes de formule générale \ce{C_{n}H_{2n+2}}.
\item La \cle{distillation fractionnée} sépare les constituants d'un mélange liquide grâce à la différence de leurs températures d'ébullition : le constituant le plus volatil (température d'ébullition la plus basse) s'échappe en haut de la colonne ; le moins volatil reste en bas.
\item Ordre des fractions, du haut vers le bas de la colonne : gaz de pétrole, essence, kérosène, gasoil, fioul, et résidu (bitume) au fond.
\item Utilisations : gaz butane (cuisson), essence (moteurs des motos et voitures), kérosène (avions, lampes), gasoil (camions, groupes électrogènes), fioul (chauffage, industrie), bitume (revêtement des routes).
\item Sécurité : les vapeurs d'hydrocarbures sont inflammables ; les hydrocarbures répandus polluent les sols, les eaux et l'air.
\end{itemize}

\courssub{Démarche et corrigé commenté}

\begin{exoresolu}[Corrigé commenté de l'activité]
\textbf{Questions 1 et 2 : courbe et évolution.} En reportant les huit points (1 ; $-162$), (4 ; $-0,5$), (5 ; 36), (6 ; 69), (8 ; 126), (10 ; 174), (12 ; 216), (16 ; 287), on obtient une courbe croissante, d'abord très pentue puis qui s'infléchit légèrement. Conclusion : \emph{la température d'ébullition d'un alcane augmente quand le nombre d'atomes de carbone de sa molécule augmente}. Les petites molécules sont les plus volatiles.
\tcblower
\textbf{Question 3 : principe de la séparation.} Le pétrole brut est un mélange d'hydrocarbures dont les températures d'ébullition sont différentes. Chauffé à environ \SI{400}{\celsius}, il se vaporise presque entièrement. Dans la colonne, la température décroît de bas en haut : chaque vapeur se condense au niveau où la température est voisine de sa température d'ébullition. On recueille ainsi des \cle{fractions} à différentes hauteurs. C'est une séparation \emph{physique} : aucune espèce chimique n'est créée.

\textbf{Question 4 : légende de la colonne.} Les molécules légères (peu de carbone, ébullition basse) montent le plus haut ; les lourdes restent en bas. D'où :
\begin{itemize}
\item Soutirage 1 (haut, le plus froid) : gaz de pétrole (propane, butane) ;
\item Soutirage 2 : essence ;
\item Soutirage 3 : kérosène ;
\item Soutirage 4 : gasoil ;
\item Soutirage 5 : fioul ;
\item Soutirage 6 et fond de colonne : résidu, le bitume.
\end{itemize}

\textbf{Question 5 : applications concrètes.} (a) Les motos-taxis roulent à l'\cle{essence}, soutirée au niveau 2. (b) Les cuisinières utilisent le \cle{gaz butane}, qui appartient aux gaz de pétrole soutirés au sommet (niveau 1) : le butane bout à \SI{-0,5}{\celsius}, il est donc gazeux à la température ambiante et se liquéfie sous pression dans la bonbonne. (c) Ces deux fractions sont donc les plus volatiles du brut.

\textbf{Question 6 : exemple d'avis citoyen.} « Monsieur le Chef, la fuite d'huile de Ndogpassi menace notre nappe phréatique et notre santé. Nous proposons trois mesures de prévention : (1) interdire les dépôts clandestins de carburant et regrouper la vente dans des stations agréées équipées de cuves étanches ; (2) récupérer les huiles usées des moteurs dans des bidons spéciaux au lieu de les verser dans les caniveaux ; (3) installer des bacs de rétention et du sable absorbant autour des points de stockage. Pour la sécurité des dépôts : (a) interdire toute flamme, cigarette ou téléphone allumé à moins de \SI{20}{m}, car les vapeurs d'essence s'enflamment facilement ; (b) stocker les bidons à l'ombre, loin des habitations, avec des extincteurs en état de marche. Comptant sur votre engagement, les élèves de la classe de 3ème. »
\end{exoresolu}

\begin{attention}[Pièges fréquents]
Ne confonds pas « le plus volatil » avec « le plus lourd » : c'est l'inverse. Le bitume, très lourd, ne s'évapore pas : il reste au fond. Autre erreur : croire que la distillation transforme chimiquement le pétrole ; elle ne fait que séparer physiquement ses constituants.
\end{attention}

\courssub{Bilan de l'activité}

Cette activité t'a permis de relier trois savoirs essentiels de la leçon. D'abord, l'exploitation du tableau de données a montré expérimentalement que la température d'ébullition des alcanes croît avec le nombre d'atomes de carbone : c'est la clé de tout le raffinage. Ensuite, le légendage de la colonne a confirmé que la distillation fractionnée classe les fractions selon leur volatilité, du gaz de pétrole au sommet jusqu'au bitume au fond. Enfin, l'avis citoyen a mis en évidence que les produits pétroliers, indispensables à la vie quotidienne camerounaise (moto-taxis, gaz domestique, groupes électrogènes), exigent une utilisation responsable : prévention des fuites, gestion des huiles usées et respect strict des consignes de sécurité. Tu as ainsi mobilisé la démarche complète : observer des données, conclure, appliquer à un schéma technique, puis agir en citoyen responsable.

\newpage

\coursec{Méthodes et exercices résolus}

\courssub{Identifier et définir les notions clés de la leçon}

\begin{methode}[Identifier et définir les notions de la leçon dans une situation concrète]
Pour répondre à une question de type « définir » ou « identifier » au BEPC :
\begin{enumerate}
\item \textbf{Repérer le mot-clé} dans l'énoncé (pétrole brut, hydrocarbure, raffinage, distillation fractionnée, coupe pétrolière...).
\item \textbf{Rappeler la définition exacte} vue en cours, avec ses mots obligatoires :
\begin{itemize}
\item \cle{Pétrole brut} : mélange naturel complexe d'hydrocarbures liquides (avec des impuretés), formé pendant des millions d'années à partir de débris d'organismes marins enfouis, et extrait des gisements souterrains ou sous-marins.
\item \cle{Hydrocarbure} : composé organique formé \emph{uniquement} d'atomes de carbone et d'hydrogène (ex. méthane \ce{CH4}, butane \ce{C4H10}).
\item \cle{Raffinage} : ensemble des opérations qui transforment le pétrole brut en produits directement utilisables.
\item \cle{Distillation fractionnée} : technique de séparation des constituants d'un mélange liquide homogène selon leurs températures d'ébullition.
\item \cle{Coupe pétrolière} (ou fraction) : produit recueilli à une hauteur donnée de la tour de distillation.
\end{itemize}
\item \textbf{Relier la définition à la situation} de l'énoncé (ex. : la SONARA de Limbé raffine le pétrole brut).
\item \textbf{Rédiger une phrase complète}, sans abréviation, en recopiant le vocabulaire scientifique correctement.
\end{enumerate}
\end{methode}

\begin{exoresolu}[Définitions et vocabulaire]
Le pétrole extrait au large de Kribi est transporté jusqu'à la raffinerie de la SONARA à Limbé. \\
1. Définis : pétrole brut ; hydrocarbure ; raffinage. \\
2. Le gaz butane des bonbonnes de cuisine est un hydrocarbure de formule \ce{C4H10}. Justifie cette affirmation.
\tcblower
\textbf{1. Définitions.}
\begin{itemize}
\item Le \textbf{pétrole brut} est un mélange naturel complexe d'hydrocarbures liquides (avec des impuretés), formé pendant des millions d'années par la décomposition d'organismes marins enfouis dans les sédiments, et extrait des gisements souterrains ou sous-marins.
\item Un \textbf{hydrocarbure} est un composé organique constitué \emph{uniquement} d'atomes de carbone et d'hydrogène.
\item Le \textbf{raffinage} est l'ensemble des opérations (essentiellement la distillation fractionnée) qui transforment le pétrole brut en produits directement utilisables : gaz, essence, gasoil, fioul, bitume...
\end{itemize}
\textbf{2. Justification.} La formule \ce{C4H10} montre que la molécule de butane contient \textbf{uniquement} des atomes de carbone (4 atomes) et d'hydrogène (10 atomes). Elle ne contient ni oxygène, ni azote, ni aucun autre élément. \\
\textbf{Conclusion :} le butane \ce{C4H10} vérifie la définition d'un hydrocarbure, donc le gaz des bonbonnes de cuisine est bien un hydrocarbure.
\end{exoresolu}

\courssub{Réaliser ou interpréter une expérience : la distillation fractionnée}

\begin{experience}[Distillation fractionnée d'un mélange au laboratoire]
\textbf{Matériel :} ballon, chauffe-ballon, colonne de fractionnement, thermomètre, réfrigérant à eau, éprouvette, mélange de deux liquides miscibles. \\
\textbf{Protocole :}
\begin{enumerate}
\item Verser le mélange dans le ballon et chauffer doucement.
\item Observer le thermomètre : il se stabilise à la température d'ébullition du liquide le plus volatil.
\item Recueillir le distillat (liquide condensé) dans l'éprouvette.
\item Poursuivre le chauffage : la température monte puis se stabilise à la température d'ébullition du second liquide, recueilli à son tour.
\end{enumerate}
\textbf{Observations :} le premier distillat est recueilli à la température d'ébullition la plus basse ; le second, à la température la plus élevée. \\
\textbf{Interprétation :} le constituant le plus volatil s'évapore en premier ; ses vapeurs se condensent dans le réfrigérant refroidi par l'eau froide. \\
\textbf{Conclusion :} la distillation fractionnée sépare les constituants d'un mélange liquide homogène grâce à leurs \cle{températures d'ébullition} différentes. C'est le principe du raffinage du pétrole.
\end{experience}

\begin{methode}[Interpréter le schéma d'une distillation fractionnée]
\begin{enumerate}
\item \textbf{Identifier les éléments} du montage : ballon + chauffage, colonne de fractionnement, thermomètre, réfrigérant (eau froide), éprouvette de recueil.
\item \textbf{Expliquer le principe} en 3 temps : le mélange est chauffé $\rightarrow$ le constituant le plus volatil (température d'ébullition la plus basse) s'évapore en premier $\rightarrow$ sa vapeur se condense dans le réfrigérant et est recueillie.
\item \textbf{Raisonner sur les températures} : dans la tour de raffinage, la température diminue du bas vers le haut ; les constituants lourds (bitume) se recueillent en bas, les légers (gaz) en haut.
\item \textbf{Nommer la grandeur de tri} : c'est toujours la \cle{température d'ébullition} qui permet la séparation.
\item \textbf{Conclure} : chaque produit recueilli s'appelle une \cle{coupe pétrolière} (ou fraction).
\end{enumerate}
\textbf{Réflexe :} « Plus un constituant est léger (molécules courtes), plus sa température d'ébullition est basse, plus il monte haut dans la tour. »
\end{methode}

\begin{exoresolu}[Distillation d'un mélange au laboratoire]
Au laboratoire du collège, on chauffe un mélange de deux liquides miscibles A et B. Le liquide A bout à \SI{65}{\celsius} et le liquide B bout à \SI{100}{\celsius}. On recueille un premier distillat lorsque le thermomètre indique \SI{65}{\celsius}. \\
1. Quel liquide est recueilli en premier ? Justifie. \\
2. Comment appelle-t-on cette technique de séparation ? \\
3. Cite deux coupes pétrolières obtenues par cette technique à la SONARA.
\tcblower
\textbf{1. Liquide recueilli en premier.} Dans une distillation, le constituant qui s'évapore en premier est celui dont la température d'ébullition est la \textbf{plus basse}. Or $\SI{65}{\celsius} < \SI{100}{\celsius}$, donc le liquide A bout avant le liquide B. Le thermomètre se stabilise à \SI{65}{\celsius} pendant que A s'évapore, puis ses vapeurs se condensent dans le réfrigérant. \\
\textbf{Conclusion :} c'est le \textbf{liquide A} qui est recueilli en premier dans le distillat. \\
\textbf{2. Nom de la technique.} Il s'agit de la \textbf{distillation fractionnée}, technique de séparation des constituants d'un mélange liquide homogène selon leurs températures d'ébullition. \\
\textbf{3. Coupes pétrolières.} À la SONARA, la distillation fractionnée du pétrole brut donne par exemple : l'\textbf{essence} (carburant des motos et voitures) et le \textbf{gasoil} (carburant des camions et des groupes électrogènes). On peut citer aussi le gaz butane, le kérosène, le fioul ou le bitume.
\end{exoresolu}

\courssub{Exploiter un schéma de tour de distillation}

\begin{popfigure}
\begin{tikzpicture}[scale=0.95, every node/.style={font=\small}]
% Four
\draw[fill=poporangeL, draw=popdark] (-4.2,0.4) rectangle (-2.2,1.6);
\node at (-3.2,1) {\textbf{Four}};
\node[font=\scriptsize] at (-3.2,0.1) {\SI{400}{\celsius}};
% Tuyau four -> tour
\draw[very thick, popdark] (-2.2,1) -- (-1,1);
% Tour
\draw[fill=popcream, draw=popdark, thick] (-1,0) rectangle (1.4,6.4);
\node[rotate=90, font=\small\bfseries] at (0.2,3.2) {Tour de distillation};
% Entrée pétrole brut
\draw[-{Stealth}, very thick, popink] (-5.6,1) -- (-4.2,1);
\node[left, font=\small] at (-5.6,1) {Pétrole brut};
% Sorties
\foreach \y/\nom/\col in {5.9/{Gaz (butane, propane)}/popgreen, 4.9/{Essence}/popblue, 3.9/{Kérosène}/poppurple, 2.9/{Gasoil}/poporange, 1.9/{Fioul}/popgold, 0.5/{Bitume}/popdark}{
  \draw[-{Stealth}, thick, \col] (1.4,\y) -- (2.6,\y);
  \node[right, font=\small] at (2.6,\y) {\nom};
}
% Flèche température
\draw[-{Stealth}, very thick, poppink] (-1.7,0.4) -- (-1.7,6);
\node[rotate=90, font=\scriptsize, poppink] at (-2.05,3.2) {température croissante};
\end{tikzpicture}
\legende{Schéma simplifié d'une tour de distillation fractionnée : le pétrole brut, chauffé à environ \SI{400}{\celsius} dans le four, entre vaporisé à la base de la tour. Les coupes légères (gaz) sont recueillies en haut, les coupes lourdes (bitume) en bas.}
\end{popfigure}

\begin{methode}[Lire et exploiter le schéma d'une tour de raffinage]
\begin{enumerate}
\item \textbf{Repérer l'entrée} : le pétrole brut est chauffé vers \SI{400}{\celsius} dans un four, puis injecté à la base de la tour.
\item \textbf{Retenir le gradient de température} : la tour est chaude en bas, de plus en plus froide vers le haut.
\item \textbf{Classer les coupes} du haut (légères) vers le bas (lourdes) : gaz (butane, propane) $\rightarrow$ essence $\rightarrow$ kérosène $\rightarrow$ gasoil $\rightarrow$ fioul $\rightarrow$ bitume.
\item \textbf{Associer chaque coupe à une utilisation} : gaz = cuisine ; essence = moteurs à essence ; kérosène = avions ; gasoil = camions, groupes électrogènes ; fioul = chauffage, navires ; bitume = revêtement des routes.
\item \textbf{Justifier toute position} par la température d'ébullition : plus la coupe est recueillie haut, plus sa température d'ébullition est faible.
\end{enumerate}
\end{methode}

\begin{exoresolu}[La tour de distillation de la SONARA]
Dans la tour de distillation fractionnée d'une raffinerie, on recueille de haut en bas : les gaz, l'essence, le kérosène, le gasoil, le fioul et le bitume. \\
1. Pourquoi le pétrole brut est-il chauffé avant d'entrer dans la tour ? \\
2. Le bitume est recueilli au bas de la tour. Que peut-on dire de sa température d'ébullition par rapport à celle de l'essence ? \\
3. Un moto-taximan fait le plein d'essence et un camionneur fait le plein de gasoil. Situe ces deux coupes dans la tour l'une par rapport à l'autre.
\tcblower
\textbf{1. Rôle du chauffage.} Le pétrole brut est chauffé à environ \SI{400}{\celsius} pour \textbf{vaporiser la quasi-totalité de ses constituants} avant l'entrée dans la tour. Ce sont ces vapeurs qui montent dans la tour et se condensent ensuite à différentes hauteurs selon leur température d'ébullition. \\
\textbf{2. Comparaison des températures d'ébullition.} La température diminue du bas vers le haut de la tour : les constituants qui se condensent en bas sont ceux dont la température d'ébullition est la \textbf{plus élevée}. Le bitume est recueilli au bas de la tour et l'essence près du sommet. \\
\textbf{Conclusion :} la température d'ébullition du bitume est \textbf{nettement supérieure} à celle de l'essence ; le bitume est un constituant lourd, l'essence un constituant léger. \\
\textbf{3. Position des deux coupes.} L'essence est recueillie \textbf{au-dessus} du gasoil dans la tour : elle est plus volatile (température d'ébullition plus faible) que le gasoil. Le gasoil, plus lourd, se condense plus bas.
\end{exoresolu}

\courssub{Réaliser un calcul simple : équilibrer une équation de combustion}

\begin{methode}[Équilibrer et exploiter une équation de combustion d'hydrocarbure]
La combustion complète d'un hydrocarbure dans le dioxygène produit \textbf{toujours} du dioxyde de carbone et de l'eau :
\[ \text{hydrocarbure} + \ce{O2} \longrightarrow \ce{CO2} + \ce{H2O} \]
\begin{enumerate}
\item \textbf{Écrire} les formules des réactifs et des produits.
\item \textbf{Équilibrer} dans l'ordre : d'abord les atomes de carbone (coefficient de \ce{CO2}), puis les atomes d'hydrogène (coefficient de \ce{H2O}), enfin les atomes d'oxygène (coefficient de \ce{O2}).
\item \textbf{Vérifier} que le nombre d'atomes de chaque élément est identique des deux côtés de la flèche.
\item \textbf{Exploiter} : compter les molécules, ou utiliser les coefficients pour des proportions simples.
\end{enumerate}
\textbf{Astuce :} si le coefficient de \ce{O2} est fractionnaire (ex. $\frac{13}{2}$), multiplier toute l'équation par 2.
\end{methode}

\begin{exoresolu}[Combustion du butane des bonbonnes de gaz]
Le gaz des bonbonnes de cuisine est le butane \ce{C4H10}. Sa combustion complète dans le dioxygène de l'air produit du dioxyde de carbone et de la vapeur d'eau. \\
1. Écris et équilibre l'équation-bilan de cette combustion. \\
2. Combien de molécules de dioxygène faut-il pour brûler complètement 2 molécules de butane ? \\
3. Pourquoi ne doit-on jamais utiliser une bonbonne de gaz dans une cuisine mal aérée ?
\tcblower
\textbf{1. Équation-bilan.}
\begin{itemize}
\item Réactifs : \ce{C4H10} et \ce{O2} ; produits : \ce{CO2} et \ce{H2O}.
\[ \ce{C4H10 + O2 -> CO2 + H2O} \]
\item Carbone : 4 atomes à gauche $\Rightarrow$ coefficient 4 devant \ce{CO2} :
\[ \ce{C4H10 + O2 -> 4CO2 + H2O} \]
\item Hydrogène : 10 atomes à gauche $\Rightarrow$ coefficient 5 devant \ce{H2O} :
\[ \ce{C4H10 + O2 -> 4CO2 + 5H2O} \]
\item Oxygène : à droite, $4 \times 2 + 5 \times 1 = 13$ atomes $\Rightarrow$ coefficient $\frac{13}{2}$ devant \ce{O2}. Pour éviter la fraction, on multiplie toute l'équation par 2 :
\[ \ce{2C4H10 + 13O2 -> 8CO2 + 10H2O} \]
\item Vérification : C : $2\times4 = 8$ et $8\times1 = 8$ ; H : $2\times10 = 20$ et $10\times2 = 20$ ; O : $13\times2 = 26$ et $8\times2 + 10\times1 = 26$. L'équation est équilibrée.
\end{itemize}
\textbf{2. Proportion.} D'après les coefficients, 2 molécules de butane réagissent avec 13 molécules de dioxygène. Il faut donc \textbf{13 molécules de \ce{O2}} pour brûler complètement 2 molécules de butane. \\
\textbf{3. Sécurité.} La combustion du butane consomme beaucoup de dioxygène. Dans une cuisine mal aérée, le dioxygène s'épuise : la combustion devient \textbf{incomplète} et produit du \textbf{monoxyde de carbone \ce{CO}}, gaz incolore, inodore et très toxique qui peut provoquer la mort. Il faut donc toujours aérer la pièce.
\end{exoresolu}

\courssub{Identifier les utilisations des produits pétroliers dans la vie courante}

\begin{methode}[Associer un produit pétrolier à son utilisation]
\begin{enumerate}
\item \textbf{Dresser mentalement le tableau} coupe $\leftrightarrow$ utilisation :
\begin{center}
\begin{tabularx}{\linewidth}{|l|X|}\hline
\rowcolor{popblueL} \textbf{Coupe pétrolière} & \textbf{Utilisation courante} \\ \hline
Gaz (butane, propane) & Bonbonnes de gaz pour la cuisine \\ \hline
Essence & Carburant des motos-taxis et voitures \\ \hline
Kérosène & Carburant des avions, lampes tempête \\ \hline
Gasoil & Camions, cars, groupes électrogènes \\ \hline
Fioul & Chauffage, navires, chaudières \\ \hline
Bitume & Revêtement (goudronnage) des routes \\ \hline
Naphta, huiles, goudrons & Matières plastiques, lubrifiants, produits chimiques \\ \hline
\end{tabularx}
\end{center}
\item \textbf{Repérer dans la situation} l'objet ou l'activité (moto-taxi, avion, route goudronnée, sachet plastique...).
\item \textbf{Identifier la coupe} correspondante et \textbf{justifier} par une propriété (volatilité, pouvoir énergétique, consistance).
\item \textbf{Citer aussi les produits dérivés} : les plastiques (sachets, bidons) sont fabriqués à partir de coupes pétrolières : c'est la \cle{pétrochimie}.
\end{enumerate}
\end{methode}

\begin{exoresolu}[Les produits pétroliers autour de nous]
Dans la ville de Douala, on observe : des motos-taxis qui font le plein à la station-service ; un avion qui décolle de l'aéroport international ; l'entreprise de travaux publics qui goudronne le boulevard ; des vendeuses qui emballent les arachides dans des sachets plastiques. \\
Pour chacune de ces quatre situations, identifie le produit pétrolier utilisé et précise s'il s'agit d'une coupe directe de la distillation ou d'un produit dérivé de la pétrochimie.
\tcblower
\begin{itemize}
\item \textbf{Moto-taxi à la station-service :} le carburant est l'\textbf{essence}, coupe légère obtenue directement par distillation fractionnée du pétrole brut.
\item \textbf{Avion qui décolle :} le carburant est le \textbf{kérosène}, coupe intermédiaire issue directement de la distillation fractionnée.
\item \textbf{Goudronnage du boulevard :} le revêtement contient du \textbf{bitume}, résidu lourd recueilli au bas de la tour de distillation.
\item \textbf{Sachets plastiques :} le plastique (polyéthylène) est un \textbf{produit dérivé de la pétrochimie} : il est fabriqué par transformation chimique de coupes pétrolières (naphta), et non recueilli directement lors de la distillation.
\end{itemize}
\textbf{Conclusion :} essence, kérosène et bitume sont des coupes directes de la distillation fractionnée ; le plastique des sachets est un produit de la pétrochimie, industrie qui transforme les coupes pétrolières en nouvelles matières.
\end{exoresolu}

\begin{savaistu}[Le pétrole au Cameroun]
Le Cameroun est un pays producteur de pétrole : les gisements sont exploités notamment au large de Kribi et dans le bassin du Rio del Rey. La \textbf{SONARA} (Société Nationale de Raffinage), située à Limbé, est la raffinerie nationale : elle transforme le pétrole brut en essence, gasoil, kérosène et gaz butane pour tout le pays. Le gaz butane des bonbonnes de cuisine est distribué par des sociétés comme SCDP, Tradex ou Ola Energy. Le pétrole est une ressource \textbf{non renouvelable} : ses réserves s'épuisent, d'où l'importance d'économiser les carburants et de développer les énergies renouvelables.
\end{savaistu}

\courssub{Appliquer les consignes de sécurité et lutter contre la pollution}

\begin{methode}[Analyser un risque lié au pétrole et proposer une mesure de lutte]
\begin{enumerate}
\item \textbf{Identifier la source de pollution} : marée noire (déversement en mer), fuite de bonbonne de gaz, fumées des pots d'échappement, sachets plastiques jetés dans la nature, dégazage des navires.
\item \textbf{Nommer les conséquences} : pollution de l'eau et des sols, mort des poissons et des oiseaux, intoxication au monoxyde de carbone, incendies et explosions, plastiques qui ne se dégradent pas et bouchent les caniveaux (inondations).
\item \textbf{Proposer des mesures de prévention} : vérifier les tuyaux et détendeurs des bonbonnes, aérer les cuisines, entretenir les moteurs, ne pas jeter les sachets dans la rue, trier et recycler les plastiques.
\item \textbf{Proposer des mesures de lutte} : barrages flottants et pompage en cas de marée noire, dispersion des nappes, recyclage, réutilisation des emballages, sensibilisation.
\item \textbf{Rédiger} la réponse en reliant clairement : source $\rightarrow$ conséquence $\rightarrow$ mesure.
\end{enumerate}
\textbf{Réflexe sécurité :} gaz inflammable + air + étincelle = explosion. On ne produit jamais d'étincelle (interrupteur, flamme) quand on sent une odeur de gaz : on ouvre les fenêtres, on ferme le robinet de la bonbonne.
\end{methode}

\begin{exoresolu}[Pollution et sécurité autour du pétrole]
\textbf{Situation A :} après de fortes pluies à Yaoundé, les caniveaux remplis de sachets plastiques débordent et provoquent des inondations. \\
\textbf{Situation B :} une ménagère sent une forte odeur de gaz dans sa cuisine le soir ; son premier réflexe est d'allumer la lumière pour vérifier la bonbonne. \\
1. Explique pourquoi les sachets plastiques provoquent ces inondations et propose deux mesures de lutte. \\
2. Le réflexe de la ménagère est-il correct ? Justifie et décris la conduite à tenir.
\tcblower
\textbf{1. Sachets plastiques et inondations.} Les plastiques sont des produits dérivés du pétrole qui se dégradent \textbf{très lentement} (plusieurs centaines d'années) dans la nature. Jetés dans les rues, ils s'accumulent et \textbf{bouchent les caniveaux} ; l'eau de pluie ne peut plus s'écouler, ce qui provoque des inondations. \\
Mesures de lutte : (a) \textbf{ne pas jeter} les sachets dans la nature mais dans des poubelles, puis trier et \textbf{recycler} les plastiques ; (b) \textbf{réutiliser} les emballages ou utiliser des paniers et sacs en tissu, et organiser des campagnes de nettoiement et de sensibilisation. \\
\textbf{2. Conduite à tenir face à une fuite de gaz.} Le réflexe d'allumer la lumière est \textbf{dangereux et incorrect} : l'interrupteur peut produire une petite \textbf{étincelle} ; or le mélange gaz butane + air est \textbf{inflammable et explosif}. L'étincelle peut déclencher une explosion. \\
Conduite à tenir :
\begin{enumerate}
\item ne toucher à \textbf{aucun appareil électrique} (ni interrupteur, ni téléphone sur place) et ne pas approcher de flamme ;
\item \textbf{fermer le robinet} de la bonbonne ;
\item \textbf{ouvrir portes et fenêtres} pour évacuer le gaz ;
\item faire appel à une personne compétente pour vérifier l'installation (tuyau, détendeur).
\end{enumerate}
\textbf{Conclusion :} face au pétrole et à ses produits, la prévention (bons gestes, tri des déchets, entretien des installations) reste le moyen de lutte le plus efficace et le moins coûteux.
\end{exoresolu}

\begin{aretenir}[L'essentiel de la leçon]
\begin{itemize}
\item Le \textbf{pétrole brut} est un mélange naturel d'\textbf{hydrocarbures} (composés formés uniquement de carbone et d'hydrogène), formé pendant des millions d'années à partir d'organismes marins.
\item Le \textbf{raffinage} repose sur la \textbf{distillation fractionnée} : séparation des constituants selon leurs \textbf{températures d'ébullition}.
\item Dans la tour, les coupes légères (gaz, essence) sortent \textbf{en haut}, les coupes lourdes (fioul, bitume) \textbf{en bas}.
\item La combustion complète d'un hydrocarbure produit \ce{CO2} et \ce{H2O} ; la combustion incomplète produit le \ce{CO}, gaz mortel.
\item Les produits pétroliers polluent (marées noires, plastiques, fumées) : prévention, tri, recyclage et gestes de sécurité sont indispensables.
\end{itemize}
\end{aretenir}

\begin{attention}[Les 5 erreurs qui coûtent des points au BEPC]
\begin{enumerate}
\item \textbf{Confondre pétrole brut et produits raffinés.} Le pétrole brut n'est pas utilisable tel quel : c'est un \emph{mélange} qu'il faut raffiner. Ne jamais écrire « on met du pétrole brut dans la moto » : on met de l'\emph{essence}, coupe issue du raffinage.
\item \textbf{Définir un hydrocarbure de façon incomplète.} Oublier le mot « uniquement » est une erreur : un hydrocarbure est formé \emph{uniquement} de carbone et d'hydrogène. L'éthanol \ce{C2H6O} n'est PAS un hydrocarbure (il contient de l'oxygène).
\item \textbf{Inverser l'ordre des coupes dans la tour.} Les gaz sortent en \emph{haut} (températures d'ébullition faibles), le bitume en \emph{bas} (température d'ébullition élevée). Réflexe : « léger en haut, lourd en bas ».
\item \textbf{Mal équilibrer la combustion.} Erreurs classiques : oublier que les produits sont \ce{CO2} et \ce{H2O}, ou ne pas vérifier le décompte des atomes d'oxygène. Toujours équilibrer dans l'ordre C, puis H, puis O, et vérifier chaque élément.
\item \textbf{Négliger les consignes de sécurité dans la rédaction.} Écrire « le gaz est dangereux » sans préciser \emph{pourquoi} (mélange gaz + air + étincelle = explosion ; combustion incomplète $\rightarrow$ \ce{CO} toxique) fait perdre des points : il faut toujours justifier le danger et citer le geste correct (fermer le robinet, aérer, aucune étincelle).
\end{enumerate}
\end{attention}

\newpage

\coursec{Exercices}

\courssub{Je vérifie mes connaissances}

\begin{exercice}[QCM : le pétrole et son raffinage]
Pour chaque affirmation, choisis la (ou les) bonne(s) réponse(s).
\begin{enumerate}
\item Le pétrole brut est :
\begin{enumerate}
\item un corps pur ;
\item un mélange d'hydrocarbures ;
\item un mélange d'eau et de gaz.
\end{enumerate}
\item L'essence est :
\begin{enumerate}
\item un corps pur ;
\item un mélange d'hydrocarbures légers ;
\item un élément chimique.
\end{enumerate}
\item Dans la colonne de distillation fractionnée, le bitume est soutiré :
\begin{enumerate}
\item en haut de la colonne ;
\item au milieu de la colonne ;
\item au bas de la colonne.
\end{enumerate}
\item La distillation fractionnée sépare les constituants du pétrole grâce à leurs différences de :
\begin{enumerate}
\item masse volumique ;
\item température d'ébullition ;
\item couleur.
\end{enumerate}
\item Le kérosène est principalement utilisé comme :
\begin{enumerate}
\item carburant pour avions ;
\item combustible pour la cuisson des aliments ;
\item revêtement des routes.
\end{enumerate}
\end{enumerate}
\end{exercice}

\begin{exercice}[Vrai ou faux, avec justification]
Réponds par \textbf{vrai} ou \textbf{faux}, puis justifie ta réponse en une ou deux phrases.
\begin{enumerate}
\item Le pétrole se forme en quelques années à partir de débris végétaux.
\item Le gaz de pétrole liquéfié (GPL) des bonbonnes de gaz domestiques est soutiré en haut de la colonne de distillation.
\item Le gasoil a une température d'ébullition plus élevée que l'essence.
\item Le raffinage du pétrole par distillation fractionnée est une transformation chimique.
\item Le Cameroun possède une raffinerie de pétrole à Limbé : la SONARA.
\end{enumerate}
\end{exercice}

\begin{exercice}[Définitions à compléter]
Recopie et complète les phrases suivantes avec les mots ou groupes de mots : \emph{hydrocarbures ; distillation fractionnée ; pétrole brut ; fraction ; combustible fossile}.
\begin{enumerate}
\item Le \ldots\ldots\ldots\ldots est un mélange complexe et visqueux, de couleur sombre, extrait du sous-sol.
\item Les \ldots\ldots\ldots\ldots sont des composés organiques formés uniquement de carbone et d'hydrogène.
\item La \ldots\ldots\ldots\ldots est l'opération qui consiste à séparer les constituants d'un mélange liquide grâce à leurs températures d'ébullition différentes.
\item Chaque produit soutiré à une hauteur donnée de la colonne de raffinage constitue une \ldots\ldots\ldots\ldots.
\item Le pétrole est un \ldots\ldots\ldots\ldots, c'est-à-dire une source d'énergie formée il y a des millions d'années et non renouvelable.
\end{enumerate}
\end{exercice}

\begin{exercice}[Combustion du butane : calcul direct]
Le butane, de formule \ce{C4H10}, est le gaz contenu dans les bonbonnes de gaz domestiques.
\begin{enumerate}
\item Recopie et équilibre l'équation de la combustion complète du butane dans le dioxygène :
\[ \ce{C4H10 + O2 -> CO2 + H2O} \]
\item Combien de moles de dioxygène sont nécessaires pour brûler complètement \num{2} moles de butane ?
\item Cite deux polluants émis lors d'une combustion \emph{incomplète} du butane (par exemple dans une cuisine mal aérée).
\end{enumerate}
\end{exercice}

\courssub{Je m'entraîne}

\begin{exercice}[Classer les fractions d'un pétrole]
On donne les températures d'ébullition de cinq hydrocarbures présents dans un pétrole brut :
\begin{center}
\begin{tabular}{|l|c|}
\hline
\rowcolor{popblueL} \textbf{Hydrocarbure} & \textbf{Température d'ébullition} \\
\hline
Butane \ce{C4H10} & \SI{-0,5}{\celsius} \\
\hline
Hexane \ce{C6H14} & \SI{69}{\celsius} \\
\hline
Dodécane \ce{C12H26} & \SI{216}{\celsius} \\
\hline
Hexadécane \ce{C16H34} & \SI{287}{\celsius} \\
\hline
Pentacosane \ce{C25H52} & \SI{402}{\celsius} \\
\hline
\end{tabular}
\end{center}
\begin{enumerate}
\item Classe ces hydrocarbures par ordre croissant de température d'ébullition.
\item En déduire l'ordre dans lequel ils sont soutirés dans la colonne de distillation, du haut vers le bas.
\item À quelle fraction commerciale associe-t-on chacun de ces hydrocarbures : gaz de pétrole, essence, kérosène, gasoil, fioul lourd ?
\item Quelle relation existe-t-il entre le nombre d'atomes de carbone d'un alcane et sa température d'ébullition ?
\end{enumerate}
\end{exercice}

\begin{exercice}[Légender une colonne de distillation]
Voici le schéma simplifié d'une colonne de distillation fractionnée :
\begin{popfigure}
\begin{center}
\begin{tikzpicture}[scale=0.9]
\draw[fill=popcream,draw=popdark,thick] (0,0) rectangle (2.4,6.4);
\draw[thick,popblue] (1.2,6.4) -- (1.2,7.2) -- (3.2,7.2);
\draw[fill=white,draw=popdark] (3.55,7.2) circle (0.35);
\node[font=\small] at (4.35,7.2) {A};
\draw[thick,popdark] (2.4,0.6) -- (4.2,0.6) node[midway,below,font=\small]{chauffage};
\draw[thick,popdark] (4.2,0.6) -- (4.2,1.6) -- (2.4,1.6);
\node[font=\small,align=center] at (4.2,0.2) {pétrole brut};
\foreach \y in {5.2,4.0,2.8,1.6}{
  \draw[thick,popblue] (2.4,\y) -- (4.0,\y);
  \draw[fill=white,draw=popdark] (4.35,\y) circle (0.35);
}
\draw[thick,popblue] (1.2,0) -- (1.2,-0.8) -- (3.2,-0.8);
\draw[fill=white,draw=popdark] (3.55,-0.8) circle (0.35);
\node[font=\small,align=center] at (-0.9,3.2) {colonne de\\ distillation};
\draw[->,thick,poporange] (5.6,1.0) -- (5.6,5.6) node[midway,right,font=\small,align=left]{température\\ décroissante};
\node[font=\small] at (4.35,5.9) {B};
\node[font=\small] at (4.35,4.7) {C};
\node[font=\small] at (4.35,3.5) {D};
\node[font=\small] at (4.35,2.3) {E};
\node[font=\small] at (4.35,-0.8) {F};
\end{tikzpicture}
\end{center}
\legende{Colonne de distillation fractionnée : les sorties sont repérées par les lettres A (sommet, gaz) à F (fond de colonne).}
\end{popfigure}
\begin{enumerate}
\item Recopie et complète le tableau suivant en attribuant à chaque lettre la fraction soutirée, choisie parmi : \emph{gaz de pétrole ; essence ; kérosène ; gasoil ; fioul ; bitume}.
\begin{center}
\begin{tabular}{|c|c|c|c|c|c|c|}
\hline
\rowcolor{popblueL} \textbf{Sortie} & A & B & C & D & E & F \\
\hline
\textbf{Fraction} & & & & & & \\
\hline
\end{tabular}
\end{center}
\item Justifie l'ordre obtenu en utilisant les températures d'ébullition des fractions.
\item Pourquoi le pétrole brut doit-il être chauffé avant d'entrer dans la colonne ?
\end{enumerate}
\end{exercice}

\begin{exercice}[L'huile de vidange du mécanicien]
À Douala, un mécanicien de motos-taxis verse chaque semaine l'huile de vidange usagée dans le caniveau derrière son atelier.
\begin{enumerate}
\item Explique deux conséquences de ce geste sur l'environnement (sol, eau, êtres vivants).
\item Propose deux solutions concrètes que ce mécanicien pourrait adopter pour gérer son huile usagée de façon responsable.
\item Rédige une courte phrase de sensibilisation (slogan) à afficher dans les ateliers de mécanique.
\end{enumerate}
\end{exercice}

\begin{exercice}[Sécurité à la station-service]
Tu es chargé de rédiger l'affiche de sécurité d'une nouvelle station-service à Bafoussam.
\begin{enumerate}
\item Rédige au moins quatre consignes de sécurité à afficher à l'entrée de la station.
\item Justifie chaque consigne par une propriété des produits pétroliers (inflammabilité, vapeurs, électricité statique, etc.).
\end{enumerate}
\end{exercice}

\courssub{Je me prépare au BEPC}

\begin{exercice}[Type BEPC : la distillation fractionnée à la SONARA]
La SONARA (Société Nationale de Raffinage), située à Limbé, traite le pétrole brut pour fournir les carburants consommés au Cameroun. Le pétrole brut est chauffé à environ \SI{400}{\celsius} puis introduit dans une colonne de distillation fractionnée. On donne :
\begin{center}
\begin{tabular}{|l|c|c|}
\hline
\rowcolor{popblueL} \textbf{Fraction} & \textbf{Intervalle d'ébullition} & \textbf{Utilisation principale} \\
\hline
Gaz de pétrole (GPL) & inférieur à \SI{40}{\celsius} & gaz domestique \\
\hline
Essence & \SI{40}{\celsius} à \SI{180}{\celsius} & carburant (moteurs à essence) \\
\hline
Kérosène & \SI{180}{\celsius} à \SI{250}{\celsius} & carburant d'aviation, éclairage \\
\hline
Gasoil & \SI{250}{\celsius} à \SI{350}{\celsius} & carburant (camions, groupes électrogènes) \\
\hline
Fioul lourd & \SI{350}{\celsius} à \SI{400}{\celsius} & combustible industriel \\
\hline
Bitume & résidu (non vaporisé) & revêtement des routes \\
\hline
\end{tabular}
\end{center}
\begin{enumerate}
\item Définis les termes suivants : \emph{pétrole brut} ; \emph{distillation fractionnée} ; \emph{fraction}.
\item Explique, en deux ou trois phrases, le principe de la séparation des constituants du pétrole dans la colonne.
\item Reproduis le schéma de la colonne de distillation et place sur ce schéma, du haut vers le bas, les six fractions du tableau.
\item Un groupe électrogène d'un quartier de Yaoundé fonctionne au gasoil. À quelle hauteur de la colonne (haut, milieu ou bas) la fraction utilisée est-elle soutirée ? Justifie.
\item La combustion complète du gasoil, assimilé à l'hexadécane \ce{C16H34}, produit du dioxyde de carbone et de l'eau. Recopie et équilibre l'équation :
\[ \ce{C16H34 + O2 -> CO2 + H2O} \]
\item Le dioxyde de carbone rejeté contribue à un phénomène climatique global. Nomme ce phénomène et cite une conséquence observable au Cameroun.
\end{enumerate}
\end{exercice}

\begin{exercice}[Type BEPC : étude de document — la production pétrolière du Cameroun]
Le graphique ci-dessous donne la production annuelle de pétrole brut du Cameroun (en millions de barils) pour quelques années :
\begin{popfigure}
\begin{center}
\begin{tikzpicture}
\begin{axis}[
  width=0.85\linewidth, height=6cm,
  xlabel={Année}, ylabel={Production (millions de barils)},
  xtick={2005,2010,2015,2020},
  xticklabels={2005,2010,2015,2020},
  ymin=0, ymax=40,
  grid=both, grid style={popline!40},
  axis line style={popdark},
  label style={font=\small}, tick label style={font=\small},
]
\addplot[popcurve,mark=*,mark size=1.5pt] coordinates {(2005,30) (2010,24) (2015,37) (2020,26)};
\end{axis}
\end{tikzpicture}
\end{center}
\legende{Production annuelle de pétrole brut du Cameroun (valeurs indicatives).}
\end{popfigure}
\begin{enumerate}
\item Lis sur le graphique la production du Cameroun en 2005, en 2015 et en 2020.
\item Calcule la variation de production entre 2005 et 2015, puis entre 2015 et 2020. Exprime chaque variation en millions de barils et précise s'il s'agit d'une hausse ou d'une baisse.
\item Calcule le pourcentage de baisse de la production entre 2015 et 2020 (arrondi à l'unité).
\item Le pétrole est dit « ressource non renouvelable ». Explique ce que cela signifie et pourquoi cela doit inciter le Cameroun à diversifier ses sources d'énergie. Cite deux sources d'énergie renouvelables exploitables au Cameroun.
\item L'exploitation pétrolière peut polluer les milieux marins (marées noires). Explique un mécanisme de cette pollution et une mesure de prévention.
\end{enumerate}
\end{exercice}

\begin{exercice}[Type BEPC : problème de synthèse — gaz domestique, sécurité et environnement]
Madame Etoundi utilise à Yaoundé une bonbonne de gaz de \SI{12,5}{kg} contenant du butane \ce{C4H10}. Un soir, elle cuisine dans une petite cuisine dont la fenêtre est fermée.
\begin{enumerate}
\item Donne la définition d'un hydrocarbure et indique les deux éléments chimiques présents dans le butane.
\item Écris et équilibre l'équation de la combustion \emph{complète} du butane.
\item La cuisine étant mal aérée, la combustion devient incomplète. Nomme le gaz toxique alors dégagé et explique pourquoi il est dangereux pour la santé.
\item Cite deux autres conséquences d'une combustion incomplète (sur les ustensiles, sur l'environnement).
\item Rédige trois conseils pratiques à donner à Madame Etoundi pour utiliser sa bonbonne de gaz en toute sécurité.
\item Le butane est obtenu lors du raffinage du pétrole. À quel niveau de la colonne de distillation fractionnée est-il soutiré ? Justifie en comparant sa température d'ébullition à celle du gasoil.
\item Le transport du pétrole par oléoduc ou par camions-citernes présente des risques de fuite. Cite deux conséquences d'une fuite de pétrole sur les terres agricoles et propose une mesure de lutte contre cette pollution.
\end{enumerate}
\end{exercice}



\newpage

\coursec{Corrigés des exercices}

\coursec{Corrigés des activités et exercices — Le pétrole}

\begin{corrige}[QCM : le pétrole et son raffinage]
\begin{enumerate}
\item \textbf{b.} Le pétrole brut est un mélange complexe d'hydrocarbures.
\item \textbf{b.} L'essence est un mélange d'hydrocarbures légers (5 à 10 atomes de carbone environ).
\item \textbf{c.} Le bitume, non vaporisé, est recueilli au bas de la colonne.
\item \textbf{b.} La séparation repose sur les températures d'ébullition différentes des constituants.
\item \textbf{a.} Le kérosène est le carburant des avions à réaction.
\end{enumerate}
\end{corrige}

\begin{corrige}[Vrai ou faux, avec justification]
\begin{enumerate}
\item \textbf{Faux.} Le pétrole se forme sur des millions d'années, à partir de débris d'organismes (surtout marins) enfouis, chauffés et comprimés dans le sous-sol.
\item \textbf{Vrai.} Le GPL (butane, propane) a les températures d'ébullition les plus basses : il se vaporise et est soutiré au sommet de la colonne.
\item \textbf{Vrai.} Le gasoil bout entre environ \SI{250}{\celsius} et \SI{350}{\celsius}, l'essence entre \SI{40}{\celsius} et \SI{180}{\celsius} : le gasoil est soutiré plus bas que l'essence.
\item \textbf{Faux.} La distillation fractionnée est une transformation \emph{physique} : on sépare les constituants sans créer de nouvelles espèces chimiques.
\item \textbf{Vrai.} La SONARA (Société Nationale de Raffinage) est installée à Limbé, dans la région du Sud-Ouest.
\end{enumerate}
\end{corrige}

\begin{corrige}[Définitions à compléter]
\begin{enumerate}
\item Le \cle{pétrole brut} est un mélange complexe et visqueux, de couleur sombre, extrait du sous-sol.
\item Les \cle{hydrocarbures} sont des composés organiques formés uniquement de carbone et d'hydrogène.
\item La \cle{distillation fractionnée} est l'opération qui consiste à séparer les constituants d'un mélange liquide grâce à leurs températures d'ébullition différentes.
\item Chaque produit soutiré à une hauteur donnée de la colonne de raffinage constitue une \cle{fraction}.
\item Le pétrole est un \cle{combustible fossile}, c'est-à-dire une source d'énergie formée il y a des millions d'années et non renouvelable.
\end{enumerate}
\end{corrige}

\begin{corrige}[Combustion du butane : calcul direct]
\begin{enumerate}
\item Équation équilibrée :
\[ \ce{2C4H10 + 13O2 -> 8CO2 + 10H2O} \]
Vérification : 8 atomes de carbone, 20 atomes d'hydrogène et 26 atomes d'oxygène de chaque côté.
\item D'après les coefficients, 2 moles de butane réagissent avec 13 moles de dioxygène. Il faut donc \num{13} moles de \ce{O2}.
\item Lors d'une combustion incomplète, il se forme du \cle{monoxyde de carbone} \ce{CO} (gaz toxique, incolore et inodore) et des \cle{suies} (particules de carbone \ce{C}).
\end{enumerate}
\end{corrige}

\begin{corrige}[Classer les fractions d'un pétrole]
\begin{enumerate}
\item Ordre croissant de température d'ébullition : butane (\SI{-0.5}{\celsius}) $<$ hexane (\SI{69}{\celsius}) $<$ dodécane (\SI{216}{\celsius}) $<$ hexadécane (\SI{287}{\celsius}) $<$ pentacosane (\SI{402}{\celsius}).
\item Plus la température d'ébullition est basse, plus le constituant monte haut dans la colonne. Du haut vers le bas : butane, hexane, dodécane, hexadécane, pentacosane.
\item Correspondances : butane $\to$ gaz de pétrole (GPL) ; hexane $\to$ essence ; dodécane $\to$ kérosène ; hexadécane $\to$ gasoil ; pentacosane $\to$ fioul lourd.
\item Plus le nombre d'atomes de carbone d'un alcane est grand, plus sa température d'ébullition est élevée.
\end{enumerate}
\end{corrige}

\begin{corrige}[Légender une colonne de distillation]
\begin{enumerate}
\item A : gaz de pétrole ; B : essence ; C : kérosène ; D : gasoil ; E : fioul ; F : bitume.
\item La température décroît du bas vers le haut de la colonne. Chaque vapeur monte jusqu'au niveau où la température est inférieure à sa température d'ébullition, puis se liquéfie et est soutirée. Les fractions légères (faible température d'ébullition) sortent donc en haut, les fractions lourdes en bas ; le bitume, non vaporisé, reste au fond.
\item Le pétrole brut est chauffé à environ \SI{400}{\celsius} pour vaporiser la quasi-totalité de ses constituants : seules des vapeurs peuvent monter dans la colonne et y être séparées.
\end{enumerate}
\end{corrige}

\begin{corrige}[L'huile de vidange du mécanicien]
\begin{enumerate}
\item Conséquences possibles : l'huile infiltre le sol et pollue durablement les terres (les plantes n'y poussent plus) ; elle est entraînée par les eaux de pluie vers les rigoles, les nappes phréatiques et les rivières, où elle forme un film qui empêche l'oxygénation de l'eau et tue poissons et organismes aquatiques.
\item Solutions : récupérer l'huile usagée dans des fûts étanches et la confier à une filière de collecte ou de recyclage (régénération) ; ne jamais la verser au sol ni dans les caniveaux, et sensibiliser les clients et apprentis.
\item Exemple de slogan : « Une huile usagée bien collectée, c'est notre eau et nos sols protégés ! »
\end{enumerate}
\end{corrige}

\begin{corrige}[Sécurité à la station-service]
Exemples de consignes et justifications :
\begin{enumerate}
\item \textbf{Interdiction de fumer ou d'utiliser une flamme} : les vapeurs d'essence sont très inflammables et peuvent s'enflammer à distance.
\item \textbf{Éteindre le moteur pendant le remplissage} : un moteur chaud ou une étincelle du démarreur peut enflammer les vapeurs de carburant.
\item \textbf{Interdiction d'utiliser le téléphone portable près des pompes} : toute source d'étincelle doit être écartée des zones où s'accumulent les vapeurs.
\item \textbf{Ne pas remplir de bidons en plastique non homologués} : l'électricité statique accumulée peut produire une étincelle ; utiliser des jerricans prévus à cet effet, posés au sol.
\item \textbf{Ne pas laisser déborder le carburant} : les produits pétroliers répandus polluent le sol et les eaux et créent un risque d'incendie.
\end{enumerate}
\end{corrige}

\begin{corrige}[Type BEPC : la distillation fractionnée à la SONARA]
\begin{enumerate}
\item \cle{Pétrole brut} : mélange naturel complexe d'hydrocarbures, liquide, visqueux et sombre, extrait du sous-sol. \cle{Distillation fractionnée} : opération physique qui sépare les constituants d'un mélange liquide grâce à leurs températures d'ébullition différentes. \cle{Fraction} : chaque produit (encore un mélange) soutiré à une hauteur donnée de la colonne.
\item Le pétrole vaporisé monte dans la colonne où la température décroît du bas vers le haut. Chaque constituant se liquéfie au niveau où la température devient inférieure à sa température d'ébullition : les fractions légères sont soutirées en haut, les fractions lourdes en bas.
\item Du haut vers le bas : gaz de pétrole (GPL), essence, kérosène, gasoil, fioul lourd, bitume (résidu au fond).
\item Le gasoil est soutiré vers le bas de la colonne (entre le milieu et le bas), car sa température d'ébullition (\SI{250}{\celsius} à \SI{350}{\celsius}) est élevée : ses vapeurs se liquéfient à un niveau encore chaud, sous le kérosène et au-dessus du fioul lourd.
\item Équation équilibrée :
\[ \ce{2C16H34 + 49O2 -> 32CO2 + 34H2O} \]
\item Le \ce{CO2} renforce l'\cle{effet de serre}, ce qui provoque le réchauffement climatique. Conséquences observables au Cameroun : saisons de pluies perturbées, sécheresses plus longues dans le Nord, érosion côtière et montée des eaux autour de Douala et Limbé.
\end{enumerate}
\end{corrige}

\begin{corrige}[Type BEPC : étude de document — la production pétrolière du Cameroun]
\begin{enumerate}
\item Lecture graphique : 2005 : environ 30 millions de barils ; 2015 : environ 37 millions de barils ; 2020 : environ 26 millions de barils.
\item Entre 2005 et 2015 : $37 - 30 = +7$ millions de barils, c'est une \textbf{hausse}. Entre 2015 et 2020 : $26 - 37 = -11$ millions de barils, c'est une \textbf{baisse}.
\item Pourcentage de baisse entre 2015 et 2020 :
\begin{align*}
\frac{37 - 26}{37} \times 100 &= \frac{11}{37} \times 100 \\
&\approx \num{29.7}\,\% \approx 30\,\%.
\end{align*}
\item « Non renouvelable » signifie que le stock de pétrole s'épuise et que sa formation demande des millions d'années : il ne se reconstitue pas à l'échelle humaine. Le Cameroun doit donc préparer l'après-pétrole. Énergies renouvelables exploitables : l'hydroélectricité (barrages sur la Sanaga), l'énergie solaire (fort ensoleillement, surtout dans le Grand Nord), éventuellement la biomasse.
\item Mécanisme : une fuite d'un oléoduc, d'un puits ou d'un navire-citerne répand le pétrole en mer ; il forme une nappe flottante qui étouffe la faune, colle aux plumes des oiseaux et pollue les côtes et les mangroves. Prévention : entretenir et contrôler régulièrement oléoducs et navires (double coque), former les équipes et préparer des plans d'intervention d'urgence (barrages flottants).
\end{enumerate}
\end{corrige}

\begin{corrige}[Type BEPC : problème de synthèse — gaz domestique, sécurité et environnement]
\begin{enumerate}
\item Un \cle{hydrocarbure} est un composé organique formé uniquement des éléments carbone et hydrogène. Le butane \ce{C4H10} contient les éléments \textbf{carbone} et \textbf{hydrogène}.
\item Combustion complète :
\[ \ce{2C4H10 + 13O2 -> 8CO2 + 10H2O} \]
\item Le gaz toxique dégagé est le \cle{monoxyde de carbone} \ce{CO}. Incolore et inodore, il se fixe sur l'hémoglobine du sang à la place du dioxygène : il empêche l'oxygénation des organes et peut provoquer maux de tête, perte de conscience, voire la mort.
\item Autres conséquences : des suies noires se déposent sur les casseroles et les murs ; le rendement de la combustion diminue (gaspillage de gaz) et les particules émises polluent l'air.
\item Conseils : cuisiner fenêtre ouverte ou dans un local ventilé ; vérifier régulièrement le tuyau et le détendeur et fermer le robinet de la bonbonne après usage ; placer la bonbonne debout, loin des flammes et des sources de chaleur, et ne jamais rechercher une fuite avec une flamme (utiliser de l'eau savonneuse).
\item Le butane est soutiré au \textbf{sommet} de la colonne : sa température d'ébullition (\SI{-0.5}{\celsius}) est très inférieure à celle du gasoil (environ \SI{250}{\celsius} à \SI{350}{\celsius}), donc il reste gazeux jusqu'en haut de la colonne.
\item Conséquences d'une fuite : le pétrole infiltre le sol et le rend stérile (les cultures ne poussent plus) ; il peut atteindre les nappes phréatiques et les cours d'eau, rendant l'eau impropre à la consommation et à l'arrosage. Mesure de lutte : entretenir et surveiller les oléoducs et camions-citernes, et en cas de fuite, confiner la zone, récupérer la terre polluée et traiter le sol (dépollution).
\end{enumerate}
\end{corrige}

\newpage

\coursec{Fiche bilan}

\begin{popfigure}
\begin{tikzpicture}[mindmap, grow cyclic, every node/.style={concept, font=\small},
root concept/.append style={concept color=popdark, font=\bfseries, text=white, minimum size=2.6cm},
level 1 concept/.append style={sibling angle=60, level distance=3.4cm, minimum size=2.1cm, font=\footnotesize\bfseries, text=white}]
\node[root concept]{Les p\'etroles}
  child[concept color=popblue]{ node {D\'efinition et origine} }
  child[concept color=popgreen]{ node {Composition : hydrocarbures} }
  child[concept color=poporange]{ node {Raffinage : distillation fractionn\'ee} }
  child[concept color=poppurple]{ node {Produits et utilisations} }
  child[concept color=poppink]{ node {Pollution et lutte} }
  child[concept color=popgold]{ node {S\'ecurit\'e et Cameroun} };
\end{tikzpicture}
\legende{Carte mentale de la le\c{c}on : les six id\'ees essentielles \`a retenir sur les p\'etroles.}
\end{popfigure}

\begin{bilan}
\textbf{D\'efinitions cl\'es}
\begin{itemize}
  \item \cle{P\'etrole brut} : m\'elange liquide naturel, noir\^atre et visqueux, form\'e pendant des millions d'ann\'ees par la d\'ecomposition d'organismes marins enfouis dans le sous-sol.
  \item \cle{Hydrocarbure} : compos\'e organique form\'e uniquement de carbone (C) et d'hydrog\`ene (H), comme le m\'ethane \ce{CH4} ou l'octane \ce{C8H18}.
  \item \cle{Raffinage} : ensemble des op\'erations qui transforment le p\'etrole brut en produits utilisables.
  \item \cle{Distillation fractionn\'ee} : s\'eparation des constituants d'un m\'elange liquide selon leurs \cle{temp\'eratures d'\'ebullition} diff\'erentes.
  \item \cle{Fraction p\'etroli\`ere} : produit obtenu dans un intervalle de temp\'erature donn\'e lors de la distillation (gaz, essence, k\'eros\`ene, gasoil, fioul, bitume).
\end{itemize}

\textbf{Les fractions de la distillation (dans l'ordre croissant des temp\'eratures d'\'ebullition)}
\begin{center}
\begin{tabularx}{\linewidth}{|l|X|X|}\hline
\rowcolor{popblueL} \textbf{Fraction} & \textbf{Temp\'erature approx.} & \textbf{Utilisation courante} \\ \hline
Gaz (butane, propane) & moins de \SI{40}{\celsius} & gaz domestique (bonbonnes) \\ \hline
Essence & \SI{40}{\celsius} \`a \SI{180}{\celsius} & carburant des motos et voitures \\ \hline
K\'eros\`ene & \SI{180}{\celsius} \`a \SI{250}{\celsius} & carburant d'avion, lampes \\ \hline
Gasoil (gazole) & \SI{250}{\celsius} \`a \SI{350}{\celsius} & camions, groupes \'electrog\`enes \\ \hline
Fioul lourd & plus de \SI{350}{\celsius} & chaudi\`eres, navires \\ \hline
Bitume (r\'esidu) & r\'esidu solide & rev\^etement des routes \\ \hline
\end{tabularx}
\end{center}

\textbf{M\'ethodes express}
\begin{itemize}
  \item \cle{Principe de la colonne \`a distiller} : le brut est chauff\'e vers \SI{400}{\celsius} ; les vapeurs montent, se refroidissent et se condensent \`a des \'etages d'autant plus hauts que la fraction est l\'eg\`ere. Les fractions l\'eg\`eres (gaz) sortent en haut, les lourdes (bitume) en bas.
  \item \cle{Combustion compl\`ete d'un hydrocarbure} : elle produit toujours du dioxyde de carbone et de l'eau : \ce{hydrocarbure + O2 -> CO2 + H2O}. Exemple : \ce{CH4 + 2O2 -> CO2 + 2H2O}.
  \item \cle{Identifier une fraction} : plus elle s'\'evapore facilement (temp\'erature d'\'ebullition basse), plus ses mol\'ecules sont petites.
  \item \cle{Lutter contre la pollution} : ne jamais jeter l'huile de vidange ou l'essence au sol ; \'eviter les feux de sachets et de pneus ; stocker les carburants dans des bidons ferm\'es, loin du feu.
\end{itemize}
\end{bilan}

\begin{aretenir}[Checklist avant le BEPC]
\begin{itemize}
  \item $\square$ Je sais d\'efinir le p\'etrole brut et raconter son origine (organismes marins, millions d'ann\'ees).
  \item $\square$ Je sais qu'un hydrocarbure ne contient que du carbone et de l'hydrog\`ene.
  \item $\square$ Je cite au moins trois exemples d'hydrocarbures : \ce{CH4}, \ce{C4H10}, \ce{C8H18}.
  \item $\square$ Je connais le principe de la distillation fractionn\'ee (s\'eparation selon les temp\'eratures d'\'ebullition).
  \item $\square$ Je sais dessiner et l\'egender une colonne de distillation (chauffage, \'etages, sorties).
  \item $\square$ Je cite les six fractions dans l'ordre : gaz, essence, k\'eros\`ene, gasoil, fioul, bitume.
  \item $\square$ J'associe chaque fraction \`a une utilisation de la vie courante (bonbonne de gaz, essence de moto, gasoil du groupe \'electrog\`ene, bitume des routes).
  \item $\square$ Je sais \'ecrire et \'equilibrer l'\'equation de combustion compl\`ete du m\'ethane.
  \item $\square$ Je cite la SONARA de Limb\'e comme raffinerie camerounaise.
  \item $\square$ Je connais les dangers des produits p\'etroliers : inflammabilit\'e, pollution de l'air, de l'eau et des sols.
  \item $\square$ Je cite au moins trois gestes de lutte contre la pollution p\'etroli\`ere.
  \item $\square$ Je respecte les consignes de s\'ecurit\'e : pas de flamme pr\`es des carburants, bidons ferm\'es et \'etiquet\'es.
\end{itemize}
\end{aretenir}

\findecours

\end{document}
